% संपूर्ण मराठी ग्रंथातील प्रकरण 74: क्रमवर्ती कलन
% हा संपादनयोग्य स्रोतखंड आहे. संकलनासाठी संपूर्ण स्रोतसंग्रहातील
% अक्षररूपे, पूर्वभाग, चिन्हव्याख्या आणि आकृती-स्रोत आवश्यक आहेत.
% मूळ: Open Logic Project; मजकूर आणि रूपांतर CC BY 4.0.
% यंत्रानुवाद, दुरुस्ती आणि तपासणी: OpenAI Codex —
% GPT-5.6 Sol आणि GPT-6 Sol, दोन्ही Ultra effort.
% दुरुस्ती-पुनरावलोकन: GPT-6.1 Sol, Ultra effort.
\chapter{क्रमवर्ती कलन}\label{pt:seq:chap}










\section{परिचय}\label{pt:seq:int:sec}

क्रमवर्ती कलन या सिद्धता पद्धती गेरहार्ड
गेन्ट्झेनने १९३० च्या दशकात मांडल्या. व्यावहारिक
सिद्ध करणे पद्धती म्हणून त्यांचा उद्देश तुलनेने कमी
होता. त्याऐवजी, सिद्धता ची संभाव्य रचना आणि
त्यांचे गुणधर्म यांविषयी काही निष्कर्ष सिद्ध करण्यासाठी
गेन्ट्झेनने त्यांचा वापर केला. नेमके असेच प्रश्न
सिद्धता-उपपत्तीच्या अभ्यासाचे विषय आहेत.

नावाप्रमाणेच क्रमवर्ती कलने सूत्रे वर
नव्हे तर \emph{क्रमवर्तींवर} काम करतात. क्रमवर्ती
म्हणजे सूत्रांच्या क्रमिकांची किंवा बहुसंचांची
जोडी. गेन्ट्झेनच्या मूळ पद्धतींमध्ये (\Log{LK} आणि
\Log{LJ}) क्रमिका वापरल्या होत्या; आता
सिद्धता-उपपत्तीचे अभ्यासक बहुसंच वापरणे पसंत करतात.
बहुसंच हा संचासारखा असतो, पण त्यात एखादा
घटक एकाहून अधिक वेळा येऊ शकतो.
संचाप्रमाणेच, मात्र, घटक चा क्रम महत्त्वाचा
नसतो. म्हणून $\{A, A, B\}$ आणि
$\{A, B, A\}$ ही एकाच बहुसंचाची रूपे आहेत.
पण हा बहुसंच $\{A,
B\}$ आणि $\{A, A, B, B\}$ यांहून वेगळा आहे,
कारण त्यांतील $A$ आणि $B$ यांच्या प्रतींच्या
संख्या वेगवेगळ्या आहेत.

सिद्धता-उपपत्तीमध्ये सूत्रांचे बहुसंच
परंपरेने मोठ्या ग्रीक अक्षरांनी लिहितात. म्हणून
$\Gamma$, $\Delta$, $\Pi$, $\Lambda$, किंवा
$\Theta$ लिहिल्यास आपण ज्या तर्कशास्त्रात काम
करतो त्यातील सूत्रांचे बहुसंच अभिप्रेत असतात.
तसेच परंपरेने, $\tuple{\Gamma, \Delta}$ असे लिहिण्याऐवजी
सिद्धता-उपपत्तीचे अभ्यासक दोन घटक क्रमवर्ती-चिन्हाने
वेगळे करतात; येथे आपण~$\Sequent$ वापरतो.

\begin{defn}[क्रमवर्ती]
\emph{क्रमवर्ती} म्हणजे पुढील रूपाची अभिव्यक्ती:
\[
\Gamma \Sequent \Delta
\]
येथे $\Gamma$ आणि $\Delta$ हे सूत्रांचे
सान्त (कदाचित रिकामे) बहुसंच असतात. $\Gamma$ ला
\emph{पूर्वांग}, तर $\Delta$ ला \emph{उत्तरांग} म्हणतात.
\end{defn}

$G$ हे~$\Gamma$ मधील सूत्रांचे संयोजन
मानू (रिकामे संयोजन~$\ltrue$); आणि $D$ हे~$\Delta$
मधील सूत्रांचे वियोजन मानू (रिकामे वियोजन
~$\lfalse$). मग क्रमवर्ती~$\Gamma \Sequent \Delta$
काही रचना~$\Struct{M}$ मध्ये सत्य असेल
तेव्हाच $G \lif D$ सत्य असते. अभिजात
तर्कशास्त्रात याचा अर्थ असा:~$\Gamma$ मधील
सूत्रांपैकी किमान एक असत्य असते किंवा~$\Delta$
मधील सूत्रांपैकी किमान एक सत्य असते.

$\Gamma$ हा सूत्रांचा बहुसंच असेल
तर~$\Gamma$ मध्ये~$A$ जोडल्यावर मिळणारा बहुसंच $\Gamma, A$
(किंवा $A, \Gamma$) असा लिहितो. $\Delta$ हाही
सूत्रांचा बहुसंच असेल तर $\Gamma, \Delta$
म्हणजे त्या दोन बहुसंचांचे बहुसंच-संयोजन: $\Gamma$
मध्ये $A$ ची \emph{प्रतींची संख्या}~n (म्हणजे~$A$
च्या $n$ प्रती) आणि $\Delta$ मध्ये ती~$m$ असेल,
तर $\Gamma, \Delta$ मध्ये $A$ ची प्रतींची संख्या
~$n+m$ असते. क्रमवर्तीच्या एका बाजूचा बहुसंच रिकामा
असेल तर बहुतेक वेळा ती जागा कोरी ठेवतो. उदाहरणार्थ,
$\Sequent A$ या क्रमवर्तीचे पूर्वांग रिकामे असून
उत्तरांगात~$A$ ची नेमकी~$1$ प्रत असते.

क्रमवर्ती कलनातील सिद्धता म्हणजे
क्रमवर्तींचा वृक्ष: त्याच्या सर्वांत वरच्या (किंवा
प्रारंभीच्या) क्रमवर्ती संबंधित पद्धतीतील स्वयंसिद्धके
असतात; आणि एखादी क्रमवर्ती एक किंवा दोन क्रमवर्तींच्या
खाली असेल, तर ती पद्धतीतील निगमन नियमाने त्यांच्यापासून
योग्य रीतीने मिळाली पाहिजे. प्रथम आपण अभिजात
तर्कशास्त्राच्या दोन भिन्न क्रमवर्ती पद्धती, $\Log{G1c}$
आणि~$\Log{G3c}$, पाहू. त्यांची स्वयंसिद्धके आणि नियम
तक्ते~\ref{pt:seq:int:tab:G1c} आणि~\ref{pt:seq:int:tab:G3c}
मध्ये दिले आहेत.








\begin{table}
\begin{tabular}{@{}c@{}c@{}}
\multicolumn{2}{@{}c@{}}{स्वयंसिद्धके}\\[2ex]
$A \Sequent A$ & $\lfalse \Sequent \quad$
\\[2ex]
\multicolumn{2}{@{}c@{}}{संरचनात्मक नियम}\\[2ex]
\Axiom$ \Gamma \fCenter \Delta $
\RightLabel{\LeftR{\Weakening}}
\UnaryInf$ A, \Gamma \fCenter \Delta$
\DisplayProof
&
\Axiom$ \Gamma \fCenter \Delta$
\RightLabel{\RightR{\Weakening}}
\UnaryInf$ \Gamma \fCenter \Delta, A$
\DisplayProof
\\[4ex]
\Axiom$ A, A, \Gamma \fCenter \Delta $
\RightLabel{\LeftR{\Contraction}}
\UnaryInf$ A, \Gamma \fCenter \Delta$
\DisplayProof
&
\Axiom$ \Gamma \fCenter \Delta, A, A$
\RightLabel{\RightR{\Contraction}}
\UnaryInf$ \Gamma \fCenter \Delta, A$
\DisplayProof
\\[4ex]
\multicolumn{2}{@{}c@{}}{तार्किक नियम}\\[2ex]
\Axiom$ \Gamma \fCenter \Delta, A $
\RightLabel{\LeftR{\lnot}}
\UnaryInf$ \lnot A, \Gamma \fCenter \Delta$
\DisplayProof
&
\Axiom$A, \Gamma \fCenter \Delta$
\RightLabel{\RightR{\lnot}}
\UnaryInf$ \Gamma \fCenter \Delta, \lnot A $
\DisplayProof
\\[4ex]\begin{tabular}[t]{@{}l@{}}
\Axiom$ A, \Gamma \fCenter \Delta$
\RightLabel{\LeftR{\land}}
\UnaryInf$ A \land B, \Gamma \fCenter \Delta$
\DisplayProof
\\[3ex]
\Axiom$B, \Gamma \fCenter \Delta$
\RightLabel{\LeftR{\land}}
\UnaryInf$A \land B, \Gamma \fCenter \Delta$
\DisplayProof\\[3ex]
\end{tabular}
&
\Axiom$\Gamma \fCenter \Delta, A$
\Axiom$ \Gamma \fCenter \Delta, B$
\RightLabel{\RightR{\land}}
\BinaryInf$ \Gamma \fCenter \Delta, A \land B $
\DisplayProof
\\[4ex]\Axiom$A, \Gamma \fCenter \Delta$
\Axiom$B, \Gamma \fCenter \Delta$
\RightLabel{\LeftR{\lor}}
\BinaryInf$A \lor B, \Gamma \fCenter \Delta$
\DisplayProof
&
\begin{tabular}[t]{@{}r@{}}
\Axiom$\Gamma \fCenter \Delta, A$
\RightLabel{\RightR{\lor}}
\UnaryInf$ \Gamma \fCenter \Delta, A \lor B$
\DisplayProof
\\[3ex]
\Axiom$ \Gamma \fCenter \Delta, B$
\RightLabel{\RightR{\lor}}
\UnaryInf$ \Gamma \fCenter \Delta, A \lor B$
\DisplayProof\\[3ex]
\end{tabular}
\\[4ex]
\Axiom$ \Gamma \fCenter \Delta, A$
\Axiom$ B, \Gamma \fCenter \Delta$
\RightLabel{\LeftR{\lif}}
\BinaryInf$ A \lif B, \Gamma \fCenter \Delta$
\DisplayProof
&
\Axiom$ A, \Gamma \fCenter \Delta, B$
\RightLabel{\RightR{\lif}}
\UnaryInf$ \Gamma \fCenter \Delta, A \lif B $
\DisplayProof
\\[4ex]
\Axiom$ A(t), \Gamma \fCenter \Delta$
\RightLabel{\LeftR{\lforall}}
\UnaryInf$ \lforall[x][A(x)],\Gamma \fCenter \Delta$
\DisplayProof
&
\Axiom$ \Gamma \fCenter \Delta, A(a) $
\RightLabel{\RightR{\lforall}}
\UnaryInf$ \Gamma \fCenter \Delta, \lforall[x][A(x)]$
\DisplayProof
\\[4ex]
\Axiom$ A(a), \Gamma \fCenter \Delta $
\RightLabel{\LeftR{\lexists}}
\UnaryInf$ \lexists[x][A(x)], \Gamma \fCenter \Delta$
\DisplayProof
&
\Axiom$ \Gamma \fCenter \Delta, A(t) $
\RightLabel{\RightR{\lexists}}
\UnaryInf$ \Gamma \fCenter \Delta, \lexists[x][A(x)]$
\DisplayProof
\end{tabular}
\caption{\Log{G1c} चे नियम. $\RightR{\forall}$ आणि
$\LeftR{\exists}$ मध्ये स्थिरांक~$a$ हा
निष्कर्ष-क्रमवर्तीमध्ये येऊ नये. $\Gamma$ आणि
$\Delta$ हे सूत्रांचे बहुसंच आहेत.}
\label{pt:seq:int:tab:G1c}
\end{table}










\begin{table}
\begin{tabular}{@{}c@{}c@{}}
\multicolumn{2}{@{}c@{}}{स्वयंसिद्धके}\\[2ex]
$C, \Gamma \Sequent \Delta, C$ & $\lfalse, \Gamma \Sequent \Delta$
\\[2ex]
\multicolumn{2}{@{}c@{}}{तार्किक नियम}\\[2ex]
\Axiom$ \Gamma \fCenter \Delta, A $
\RightLabel{\LeftR{\lnot}}
\UnaryInf$ \lnot A, \Gamma \fCenter \Delta$
\DisplayProof
&
\Axiom$A, \Gamma \fCenter \Delta$
\RightLabel{\RightR{\lnot}}
\UnaryInf$ \Gamma \fCenter \Delta, \lnot A $
\DisplayProof
\\[4ex]
\Axiom$ A, B, \Gamma \fCenter \Delta$
\RightLabel{\LeftR{\land}}
\UnaryInf$ A \land B, \Gamma \fCenter \Delta$
\DisplayProof
&
\Axiom$\Gamma \fCenter \Delta, A$
\Axiom$ \Gamma \fCenter \Delta, B$
\RightLabel{\RightR{\land}}
\BinaryInf$ \Gamma \fCenter \Delta, A \land B $
\DisplayProof
\\[4ex]\Axiom$A, \Gamma \fCenter \Delta$
\Axiom$B, \Gamma \fCenter \Delta$
\RightLabel{\LeftR{\lor}}
\BinaryInf$A \lor B, \Gamma \fCenter \Delta$
\DisplayProof
&
\Axiom$\Gamma \fCenter \Delta, A, B$
\RightLabel{\RightR{\lor}}
\UnaryInf$ \Gamma \fCenter \Delta, A \lor B$
\DisplayProof
\\[4ex]
\Axiom$ \Gamma \fCenter \Delta, A$
\Axiom$ B, \Gamma \fCenter \Delta$
\RightLabel{\LeftR{\lif}}
\BinaryInf$ A \lif B, \Gamma \fCenter \Delta$
\DisplayProof
&
\Axiom$ A, \Gamma \fCenter \Delta, B$
\RightLabel{\RightR{\lif}}
\UnaryInf$ \Gamma \fCenter \Delta, A \lif B $
\DisplayProof
\\[4ex]
\Axiom$ A(t), \lforall[x][A(x)], \Gamma \fCenter \Delta$
\RightLabel{\LeftR{\lforall}}
\UnaryInf$ \lforall[x][A(x)],\Gamma \fCenter \Delta$
\DisplayProof
&
\Axiom$ \Gamma \fCenter \Delta, A(a) $
\RightLabel{\RightR{\lforall}}
\UnaryInf$ \Gamma \fCenter \Delta, \lforall[x][A(x)]$
\DisplayProof
\\[4ex]
\Axiom$ A(a), \Gamma \fCenter \Delta $
\RightLabel{\LeftR{\lexists}}
\UnaryInf$ \lexists[x][A(x)], \Gamma \fCenter \Delta$
\DisplayProof
&
\Axiom$ \Gamma \fCenter \Delta, \lexists[x][A(x)], A(t) $
\RightLabel{\RightR{\lexists}}
\UnaryInf$ \Gamma \fCenter \Delta, \lexists[x][A(x)]$
\DisplayProof
\end{tabular}
\caption{\Log{G3c} चे नियम. $\RightR{\forall}$ आणि
$\LeftR{\exists}$ मध्ये स्थिरांक~$a$ हा
निष्कर्ष-क्रमवर्तीमध्ये येऊ नये. $\Gamma$ आणि
$\Delta$ हे सूत्रांचे बहुसंच आहेत.
स्वयंसिद्धकांमध्ये सूत्र~$C$ अण्वीय असले पाहिजे.}
\label{pt:seq:int:tab:G3c}
\end{table}




या पद्धतींमध्ये सूक्ष्म फरक आहेत आणि त्यामुळे
त्यांचे सिद्धता-उपपत्तीय गुणधर्म वेगवेगळे आहेत.
\Log{G1c} ची स्वयंसिद्धके $A
\Sequent A$ अशी असतात; त्यात इतर कोणतेही
सूत्रांचालत नाहीत. \Log{G3c} ची स्वयंसिद्धके
$C, \Gamma \Sequent \Delta, C$ या रूपाची
असतात; त्यात~$\Gamma$ आणि~$\Delta$ मधील इतर
``बाजूची'' सूत्रे मनमानी असू शकतात, पण
दोन्ही बाजूंना येणारे~$C$ हे सूत्र
\emph{अण्वीय} असले पाहिजे. \Log{G1c} मध्ये
$\LeftR{\land}$ आणि $\RightR{\lor}$ नियमांच्या
प्रत्येकी दोन आवृत्त्या आहेत, तर \Log{G3c} मध्ये
प्रत्येकी एकच आहे. तसेच \Log{G1c} मध्ये
``आकुंचन'' (\LeftR{\Contraction}, \RightR{\Contraction})
आणि ``दुर्बलीकरण'' (\LeftR{\Weakening}, \RightR{\Weakening})
हे अतिरिक्त ``संरचनात्मक नियम'' आहेत.
यांशिवाय असत्याचे स्वयंसिद्धकही आहे: \Log{G1c} मध्ये
$\lfalse\Sequent\quad$, आणि \Log{G3c} मध्ये
$\lfalse,\Gamma\Sequent\Delta$. त्यामुळे समान
अणुरूप सूत्र दोन्ही बाजूंना असण्याचा प्रसंग हाच
स्वयंसिद्धकांचा एकमेव प्रसंग नाही.

\Log{G1c} आणि \Log{G3c} या दोन्ही
पद्धती अभिजात तर्कशास्त्रासाठी निर्दोष व संपूर्ण
आहेत. म्हणजे या पद्धतींमध्ये सिद्धतायोग्य असलेली
प्रत्येक क्रमवर्ती प्रत्येक संरचनांमध्ये सत्य
असते; आणि प्रत्येक संरचनांमध्ये सत्य असलेल्या
प्रत्येक क्रमवर्तीची सिद्धता असते. म्हणून एखाद्या
क्रमवर्तीची सिद्धता \emph{आहे का}, या प्रश्नाचे
उत्तर तर्कशास्त्राच्या इतर पद्धतींनीही (उदा.,
प्रतिमान-उपपत्तीने) मिळवता येते. आणि दोन्ही पद्धती
नेमक्या प्रत्येक संरचनांमध्ये सत्य असलेल्या
क्रमवर्ती सिद्ध करतात; म्हणून त्या विशेषतः
त्याच क्रमवर्ती सिद्ध करतात.

पण सिद्धता-उपपत्तीला एखादी गोष्ट
\emph{सिद्धतायोग्य आहे का} एवढाच नव्हे तर ती
\emph{कशी} सिद्ध होते हाही प्रश्न महत्त्वाचा असतो.
सिद्धता-उपपत्तीचे अभ्यासक असे प्रश्न विचारतात:
``एखादा विशिष्ट नियम नेहमी टाळता येईल का?'',
``नव्या क्रमवर्ती सिद्धतायोग्य न होता हा नियम पद्धतीत
जोडता येईल का?'', किंवा ``एका पद्धतीतील सिद्धता
दुसऱ्या पद्धतीतील सिद्धता मध्ये रूपांतरित करता
येतील का?'' अशा प्रश्नाचे उत्तर ``होय'' असेल, तर
त्या तथ्याची सिद्धता बहुधा सिद्धता च्या
गुंतागुंतीवर किंवा, समतुल्यपणे, त्या सिद्धता च्या रचनेवर
विगमन करून दिली जाते. त्यातून केवळ तथ्याचीच
सिद्धता मिळत नाही (उदा., \Log{G1c} ने एखादी
क्रमवर्ती सिद्ध केली तर तीच क्रमवर्ती \Log{G3c}
नेही सिद्ध होते), तर एक प्रक्रिया मिळते (उदा.,
\Log{G1c} मधील सिद्धता ना \Log{G3c} मधील
सिद्धता मध्ये रूपांतरित करण्याची).

सिद्धता-उपपत्तीतील एक मध्यवर्ती निष्कर्ष म्हणजे
\emph{कट-निरसन प्रमेय}. त्यानुसार पुढील कट नियम
\[
\Axiom$\Gamma \fCenter \Delta, A$
\Axiom$A, \Pi \fCenter \Lambda$
\RightLabel{\Cut}
\BinaryInf$\Gamma, \Pi \fCenter \Delta, \Lambda$
\DisplayProof
\]
क्रमवर्ती पद्धतींमध्ये जोडला तरी कोणत्या क्रमवर्ती
सिद्धतायोग्य आहेत यांत बदल होत नाही. याहूनही,
त्या निष्कर्षाची सिद्धता अशी पद्धत देते की
\Log{G1c} (किंवा \Log{G3c}) चे नियम आणि~\Cut{}
वापरून केलेल्या कोणत्याही सिद्धता चे रूपांतर
फक्त \Log{G1c} (किंवा \Log{G3c}) चेच नियम
वापरणाऱ्या सिद्धतेत करता येते. दुसऱ्या शब्दांत, ही
प्रक्रिया~\Cut{} वापरणाऱ्या सिद्धता मधून
हा नियम ``निरसित'' करते; म्हणून हे नाव.













\section{नियम आणि सिद्धता}\label{pt:seq:rp:sec}

आधी काही व्याख्या देऊ आणि काही संज्ञांचा अर्थ निश्चित करू.

उदाहरणार्थ, $A$ आणि~$B$ असलेल्या क्रमवर्तींपासून
$A \land B$ असलेल्या क्रमवर्ती सिद्ध करणे करू देणारे \Log{G3c} मधील
दोन नियम पाहा:
\[
\Axiom$ A, B, \Gamma \fCenter \Delta$
\RightLabel{\LeftR{\land}}
\UnaryInf$ A \land B, \Gamma \fCenter \Delta$
\DisplayProof
\qquad
\Axiom$\Gamma \fCenter \Delta, A$
\Axiom$ \Gamma \fCenter \Delta, B$
\RightLabel{\RightR{\land}}
\BinaryInf$ \Gamma \fCenter \Delta, A \land B $
\DisplayProof
\]
रेषेच्या वरच्या क्रमवर्तींना \emph{आधारविधाने}, तर खालच्या
क्रमवर्तीला \emph{निष्कर्ष} म्हणतात. प्रत्येक नियमातील
$\Gamma$ आणि $\Delta$ यांतील सूत्रे ना \emph{संदर्भ}
किंवा \emph{बाजूची सूत्रे} म्हणतात. निष्कर्षातील
संदर्भाबाहेरचे सूत्र हे \emph{मुख्य} सूत्र असते;
आणि आधारविधानांतील संदर्भाबाहेरच्या सूत्रे ना
\emph{सक्रिय} (किंवा \emph{सहायक}) सूत्रे म्हणतात.

परिमाणकांचे नियम थोडे अधिक गुंतागुंतीचे आहेत. उदाहरणार्थ,
$\lforall$ साठीचे \Log{G1c} मधील नियम असे आहेत:
\[\Axiom$ A(t), \Gamma \fCenter \Delta$
\RightLabel{\LeftR{\lforall}}
\UnaryInf$ \lforall[x][A(x)],\Gamma \fCenter \Delta$
\DisplayProof
\qquad
\Axiom$ \Gamma \fCenter \Delta, A(a) $
\RightLabel{\RightR{\lforall}}
\UnaryInf$ \Gamma \fCenter \Delta, \lforall[x][A(x)]$
\DisplayProof
\]
\LeftR{\lforall} नियमातील सक्रिय सूत्र हे~$A(t)$ आहे;
येथे $t$ हे बंद पद, म्हणजे चर नसलेले पद, आहे. एखादे
निगमन योग्य असते—म्हणजे ते \LeftR{\lforall} या
रूपरेषात्मक नियमाशी जुळते—जर एखादे सूत्र~$A$,
चर~$x$ आणि बंद पद~$t$ असे असतील की निष्कर्षातील मुख्य
सूत्र $\lforall[x][A]$ असेल आणि आधारविधानातील
सक्रिय सूत्र हे~$\Subst{A}{t}{x}$ असेल.
\RightR{\lforall} नियमही याचप्रमाणे कार्य करतो; मात्र
कोणत्याही $t$ पदाऐवजी आधारविधानातील सक्रिय सूत्र
$\Subst{A}{a}{x}$ असे असले पाहिजे, जेथे $a$ हा
स्थिरांक आहे. या अचराला \RightR{\lforall}
निगमनाचा \emph{आयगेन चर} म्हणतात. खालच्या क्रमवर्तीमध्ये
$a$ येत नसेल—म्हणजेच $\Gamma$, $\Delta$ किंवा~$A$
यांत $a$ नसेल—तरच हे निगमन योग्य असते. ही
\emph{आयगेन चराची अट} होय.

\Log{G1c} आणि \Log{G3c} यांसारख्या क्रमवर्ती पद्धतींत अशा
रूपरेषात्मक रीतीने दिलेल्या नियमांचा संग्रह आणि स्वयंसिद्धके
म्हणून ग्राह्य अशा, त्याच रीतीने दिलेल्या, काही क्रमवर्ती असतात.
अशा पद्धतीतील सिद्धता म्हणजे स्वयंसिद्धकांपासून निगमन नियमांनी
विगमनाने तयार होणारे क्रमवर्तींचे सान्त वृक्ष. प्रत्येक पान
स्वयंसिद्धक असते; आणि इतर प्रत्येक क्रमवर्ती तिच्या लगेच
वरच्या क्रमवर्तींपासून पद्धतीतील एका निगमन नियमाने मिळते.

\begin{defn}\label{pt:seq:rp:defn:proof}
क्रमवर्ती पद्धतीतील \emph{सिद्धता} म्हणजे खालीलप्रमाणे
विगमनाने बनवलेले क्रमवर्तींचे सान्त वृक्ष:
\begin{enumerate}
  \item\label{pt:seq:rp:defn:prf:axiom} प्रत्येक स्वयंसिद्धक ही सिद्धता असते.
  \item\label{pt:seq:rp:defn:prf:one-prem} $\pi_1$ ही $S_1$ या क्रमवर्तीवर
  संपणारी सिद्धता असेल आणि $S_1$ पासून एका आधारविधानाच्या
  $R$ नियमाने $S$ ही क्रमवर्ती मिळत असेल, तर
  \[
  \AxiomC{}
  \RightLabel{$\pi_1$}
  \DeduceC{$S_1$}
  \RightLabel{$R$}
  \UnaryInfC{$S$}
  \DisplayProof
  \]
  ही सिद्धता असते.
  \item\label{pt:seq:rp:defn:prf:two-prem} $\pi_1$ आणि $\pi_2$ या
  अनुक्रमे $S_1$ आणि~$S_2$ या क्रमवर्तींवर संपणाऱ्या सिद्धता
  असतील, आणि $S_1$ व $S_2$ यांच्यापासून दोन आधारविधानांच्या $R$ नियमाने
  $S$ मिळत असेल, तर
  \[
  \AxiomC{}
  \RightLabel{$\pi_1$}
  \DeduceC{$S_1$}
  \AxiomC{}
  \RightLabel{$\pi_2$}
  \DeduceC{$S_2$}
  \RightLabel{$R$}
  \BinaryInfC{$S$}
  \DisplayProof
  \]
  ही सिद्धता असते.
\end{enumerate}
\end{defn}

\ref{pt:seq:rp:defn:proof} मध्ये सिद्धता ची व्याख्या क्रमवर्तींचे वृक्ष म्हणून
दिली आहे. हे वृक्ष ``वरच्या दिशेने'' वाढतात असे आपण मानतो.
सर्वांत वरची (पानांवरील) स्थाने स्वयंसिद्धके असतात; त्यांना
\emph{आरंभीच्या क्रमवर्ती} म्हणतात. सर्वांत खालची क्रमवर्ती
\emph{अंतिम क्रमवर्ती} असते. $\pi$ ही $S$ अंतिम क्रमवर्ती
असलेली सिद्धता असेल तर $\pi \Proves S$ असेही लिहितो.
$S$ ची सिद्धता असेल तर $\Proves S$ असे लिहितो; काही वेळा
ज्या पद्धतीत सिद्धता आहे तीही $\Proves$ च्या डावीकडे दाखवतो,
उदा., $\Log{G1c} \Proves A, B \Sequent A \land
B$. आरंभीच्या क्रमवर्ती वगळता सिद्धता मधील प्रत्येक
क्रमवर्ती एखाद्या $R$ नियमाने झालेल्या निगमनाचा
\emph{निष्कर्ष} असते. सिद्धता मधील एखाद्या क्रमवर्तीच्या
लगेच वर असलेल्या एक किंवा दोन क्रमवर्तींना (वृक्षातील तिच्या
पालक-स्थानांना) त्या निगमनाची \emph{आधारविधाने} म्हणतात.

एखादी सिद्धता विगमनाने बनवताना वापरलेल्या सिद्धता ना तिच्या
\emph{उप-सिद्धता} म्हणतात. एखाद्या निगमनाच्या आधारविधानांवर
संपणाऱ्या उप-सिद्धता ना त्या निगमनाच्या \emph{लगेचच्या
उप-सिद्धता} म्हणतात.

अनेकदा सिद्धता ची गुंतागुंत एखाद्या नैसर्गिक संख्येने
मोजावी लागते. सिद्धता मधील क्रमवर्तींची संख्या मोजता येईल;
पण अनेकदा अधिक उपयोगी मोजमाप म्हणजे सिद्धता ची उंची:
अंतिम क्रमवर्तीपासून आरंभीच्या क्रमवर्तीपर्यंतच्या सर्वाधिक
निगमन-पायऱ्यांची संख्या. तिची विगमनात्मक व्याख्याही देता येते.

\begin{defn}
सिद्धता~$\pi$ ची \emph{उंची}~$\pheight{\pi}$ पुढीलप्रमाणे
विगमनाने ठरते.
\begin{enumerate}
  \item $\pi$ मध्ये एकच स्वयंसिद्धक असेल, तर $\pheight{\pi} = 0$.
  \item $\pi$ चा शेवट एका आधारविधानाच्या निगमनाने होत असेल आणि
  तिची लगेचची उप-सिद्धता~$\pi_1$ असेल, तर $\pheight{\pi} = \pheight{\pi_1} + 1$.
  \item $\pi$ चा शेवट दोन आधारविधानांच्या निगमनाने होत असेल आणि
  तिच्या लगेचच्या उप-सिद्धता~$\pi_1$ आणि~$\pi_2$ असतील, तर $\pheight{\pi} =
  \max(\pheight{\pi_1}, \pheight{\pi_2}) + 1$.
\end{enumerate}
$S$ या क्रमवर्तीची $\le n$ उंचीची सिद्धता असेल, तर $\Proves[n] S$ असे लिहितो.
\end{defn}

\begin{ex}
\Log{G3c} मध्ये $C, \Gamma \Sequent \Delta, C$ या रूपाची
प्रत्येक क्रमवर्ती स्वयंसिद्धक असते; येथे $C$ अण्वीय असले पाहिजे.
$D$ आणि $E$ ही अण्वीय वाक्ये असोत. मग $D, E \Sequent D$
आणि $D, E \Sequent E$ ही $C, \Gamma \Sequent \Delta, C$
या स्वयंसिद्धकाची उदाहरणे आहेत. उदाहरणार्थ, $C \ident E$,
$\Gamma = \{D\}$ आणि $\Delta = \emptyset$ घेतल्यास
$C, \Gamma \Sequent \Delta, C$ पासून नेमकी $D, E
\Sequent E$ ही क्रमवर्ती मिळते. (लक्षात घ्या, आपण \emph{बहुसंच} वापरतो;
म्हणून $D, E \Sequent E$ आणि $E, D \Sequent E$ या एकाच
क्रमवर्ती आहेत.)

$D, E \Sequent D$ ही \Log{G3c} च्या स्वयंसिद्धकाचे उदाहरण
असल्यामुळे ती स्वतःच \Log{G3c} मधील सिद्धता असते. त्याचप्रमाणे
$D, E \Sequent E$ हीही \ref{pt:seq:rp:defn:proof}\ref{pt:seq:rp:defn:prf:axiom}
नुसार तशीच असते. त्यांना अनुक्रमे $\pi_1$ आणि~$\pi_2$ म्हणू.
\ref{pt:seq:rp:defn:prf:two-prem} नुसार, या दोन सिद्धता घेऊन
दोन आधारविधानांच्या $\RightR{\land}$ नियमाने नवी सिद्धता
($\pi_3$) तयार करता येते:
\[
\Axiom$D, E \fCenter D$
\Axiom$D, E \fCenter E$
\RightLabel{\RightR{\land}}
\BinaryInf$D, E  \fCenter D \land E $
\DisplayProof
\]
$\pi_3$ पासून पुढे सिद्धता~$\pi_4$ मिळते:
\[
\Axiom$D, E \fCenter D$
\Axiom$D, E \fCenter E$
\RightLabel{\RightR{\land}}
\insertBetweenHyps{\hspace{5ex}}
\BinaryInf$D, E  \fCenter D \land E $
\LeftSubproofLabel{$\pi_3$}
\RightLabel{\LeftR{\land}}
\UnaryInf$D \land E  \fCenter D \land E$
\RightSubproofLabel{$\pi_4$}
\DisplayProof
\]
येथे एका आधारविधानाचा $\LeftR{\land}$ नियम वापरला आहे
(\ref{pt:seq:rp:defn:proof}\ref{pt:seq:rp:defn:prf:one-prem} नुसार). $\pi_4$ मधील
शेवटच्या निगमनाची लगेचची उप-सिद्धता ही~$\pi_3$ आहे.
$\RightR{\land}$ निगमनाच्या लगेचच्या उप-सिद्धता $\pi_1$
आणि~$\pi_2$ आहेत. $\pi_4$ च्या उप-सिद्धता मध्ये $\pi_1$,
$\pi_2$, $\pi_3$ आणि स्वतः $\pi_4$ यांचा समावेश होतो.
$\pheight{\pi_1} = \pheight{\pi_2} = 0$,
$\pheight{\pi_3} = 1$ आणि $\pheight{\pi_4} = 2$.
म्हणून कोणत्याही अण्वीय $D$ आणि~$E$ साठी
$\Log{G3c} \Proves[2] D \land E  \fCenter D
\land E$ हे दाखवले आहे.
\end{ex}













\section{सिद्धतांची उदाहरणे}\label{pt:seq:ex:sec}

क्रमवर्तींच्या सिद्धता तयार करताना बहुतेक
वेळा अंतिम क्रमवर्तीपासून वरच्या दिशेने काम करणे
सोयीचे ठरते: त्यातील एखादे सूत्र सक्रिय
सूत्र म्हणून निवडायचे, त्या क्रमवर्तीच्या वर
संबंधित नियमाचे पूर्वपक्ष लिहायचे, आणि स्वयंसिद्धके
मिळेपर्यंत ही प्रक्रिया पुन्हा करायची. उदाहरणार्थ,
पुढील क्रमवर्ती सिद्ध करायची आहे असे माना:
\[(C \land D) \lif E \Sequent C \lif (D \lif E)\]
$C \lif (D \lif E)$ विचारात घ्या: ते $A \lif B$
या रूपाचे असून क्रमवर्तीच्या उजवीकडे आहे. पूर्ण
क्रमवर्तीची \Log{G1c} च्या $\RightR{\lif}$ नियमाशी
जुळवणी करू:
\[\Axiom$ A, \Gamma \fCenter \Delta, B$
\RightLabel{\RightR{\lif}}
\UnaryInf$ \Gamma \fCenter \Delta, A \lif B $
\DisplayProof\]
यावरून $\Gamma = \{(C \land D) \lif E\}$,
$\Delta =
\emptyset$, $A \ident C$ आणि $B \ident D \lif E$
असे घ्यावे लागते. मग सिद्धतेचे शेवटचे निगमन असे असू शकेल:
\[\Axiom$ C, (C \land D) \lif E \fCenter D \lif E$
\RightLabel{\RightR{\lif}}
\UnaryInf$(C \land D) \lif E \fCenter C \lif (D \lif E)$
\DisplayProof\]
हीच कल्पना पुन्हा लावून $D \lif E$ ला सक्रिय
सूत्र मानल्यास पुढील मिळते:
\[\Axiom$ D, C, (C \land D) \lif E \fCenter E$
\RightLabel{\RightR{\lif}}
\UnaryInf$C, (C \land D) \lif E \fCenter D \lif E$
\RightLabel{\RightR{\lif}}
\UnaryInf$(C \land D) \lif E \fCenter C \lif (D \lif E)$
\DisplayProof\]
आता $(C \land D) \lif E$ वर लक्ष केंद्रित करून
\Log{G1c} चा $\LeftR{\lif}$ नियम वापरावा लागेल:
\[\Axiom$ \Gamma \fCenter \Delta, A$
\Axiom$ B, \Gamma \fCenter \Delta$
\RightLabel{\LeftR{\lif}}
\BinaryInf$ A \lif B, \Gamma \fCenter \Delta$
\DisplayProof\]
त्यातून पुढील मिळते:
\[\Axiom$D, C \fCenter E, C \land D$
\Axiom$D, C, E \fCenter E$
\RightLabel{\LeftR{\lif}}
\BinaryInf$ D, C, (C \land D) \lif E \fCenter E$
\RightLabel{\RightR{\lif}}
\UnaryInf$C, (C \land D) \lif E \fCenter D \lif E$
\RightLabel{\RightR{\lif}}
\UnaryInf$(C \land D) \lif E \fCenter C \lif (D \lif E)$
\DisplayProof\]
वरच्या डाव्या क्रमवर्तीतील उजवीकडचे सूत्र
$C \land D$ प्रमुख मानून $\RightR{\land}$ नियम
लावल्यास पुढील मिळते:
\[
\Axiom$D, C \fCenter E, C$
\Axiom$D, C \fCenter E, D$
\RightLabel{\RightR{\land}}
\BinaryInf$D, C \fCenter E, C \land D$
\Axiom$D, C, E \fCenter E$
\RightLabel{\LeftR{\lif}}
\BinaryInf$ D, C, (C \land D) \lif E \fCenter E$
\RightLabel{\RightR{\lif}}
\UnaryInf$C, (C \land D) \lif E \fCenter D \lif E$
\RightLabel{\RightR{\lif}}
\UnaryInf$(C \land D) \lif E \fCenter C \lif (D \lif E)$
\DisplayProof
\]
आतापर्यंत सगळे ठीक आहे. आपण वापरलेले नियम
\Log{G1c} आणि~\Log{G3c} मध्ये सारखे असल्याने हे
सर्व \Log{G3c} मध्येही चालते. या सिद्धता च्या
अपूर्ण भागातील वरच्या सर्व क्रमवर्तींच्या डावीकडे
आणि उजवीकडे तेच सूत्र येते; पण त्या अद्याप
अगदी स्वयंसिद्धके झालेल्या नाहीत.

\Log{G1c} मध्ये वरच्या क्रमवर्तींच्या
सिद्धता $A \Sequent A$ या रूपाच्या स्वयंसिद्धकांपासून
द्याव्या लागतात. दुर्बलीकरण नियमांनी ते सहज शक्य
आहे. उदाहरणार्थ, सर्वांत डावीकडची वरची क्रमवर्ती
$D, C \Sequent
E, C$ पुढीलप्रमाणे सिद्ध करता येते:
\[
\Axiom$C \fCenter C$
\RightLabel{\LeftR{\Weakening}}
\UnaryInf$D, C \fCenter C$
\RightLabel{\RightR{\Weakening}}
\UnaryInf$D, C \fCenter E, C$
\]
एकत्रितपणे, \Log{G1c} मधील सिद्धता अशी दिसते:
\[
\Axiom$C \fCenter C$
\RightLabel{\LeftR{\Weakening}}
\UnaryInf$D, C \fCenter C$
\RightLabel{\RightR{\Weakening}}
\UnaryInf$D, C \fCenter E, C$
\Axiom$D \fCenter D$
\RightLabel{\LeftR{\Weakening}}
\UnaryInf$D, C \fCenter D$
\RightLabel{\RightR{\Weakening}}
\UnaryInf$D, C \fCenter E, D$
\RightLabel{\RightR{\land}}
\BinaryInf$D, C \fCenter E, C \land D$
\Axiom$E \fCenter E$
\RightLabel{\LeftR{\Weakening}}
\UnaryInf$C, E \fCenter E$
\RightLabel{\LeftR{\Weakening}}
\UnaryInf$D, C, E \fCenter E$
\RightLabel{\LeftR{\lif}}
\BinaryInf$ D, C, (C \land D) \lif E \fCenter E$
\RightLabel{\RightR{\lif}}
\UnaryInf$C, (C \land D) \lif E \fCenter D \lif E$
\RightLabel{\RightR{\lif}}
\UnaryInf$(C \land D) \lif E \fCenter C \lif (D \lif E)$
\DisplayProof
\]
या सिद्धतेची रचना (विशेषतः तिची उंची) सूत्रे
$C$, $D$, आणि~$E$ काय आहेत यावर अवलंबून नाही.
वरील सिद्धता मध्ये या चिन्हांच्या जागी कोणतीही
सूत्रे घातली तरी मिळणारी रचना~\Log{G1c}
मधील सिद्धता राहते. \Log{G1c} मधील ही
सिद्धता \emph{रूपरेषात्मक} आहे.

\Log{G3c} मध्ये स्वयंसिद्धके
$A, \Gamma \Sequent
\Delta, A$ या रूपाची क्रमवर्ती असतात; त्यात~$A$
अण्वीय असले पाहिजे. म्हणून $D, C \Sequent
E, C$ ही \Log{G3c} ची स्वयंसिद्धकासारखी दिसत
असली ($\Gamma = \{D\}$; $\Delta =
\{E\}$), तरी~$C$ प्रत्यक्षात अण्वीय असेल तेव्हाच
ती \emph{खरोखर} स्वयंसिद्धक असते. $C$, $D$, आणि $E$ ही अण्वीय सूत्रे
असतील, तर हा सिद्धता-वृक्ष \Log{G3c} मध्ये पूर्ण
सिद्धता असतो. ही पुरेशी अट आहे; इतर बाजूच्या सूत्रामुळे
एखादी वरची क्रमवर्ती स्वयंसिद्धक असण्याची शक्यताही आहे.

काटेकोरपणे पाहता, आपण अजून
\[\Log{G3c} \Proves (C \land D) \lif E \fCenter C \lif (D \lif
E),\] हे दाखवलेले नाही. तरी व्यवहारात एवढेच
करणे आपल्याला पुरेसे आहे (आणि शक्यही आहे).
कारण $A, \Gamma \Sequent \Delta, A$ या रूपाची
प्रत्येक क्रमवर्ती \Log{G3c} मध्ये सिद्धतायोग्य
असते, जरी~$A$ स्वतः अण्वीय नसले तरी.
हे~$A$ च्या खोली~$\depth{A}$ वर विगमन करून
दाखवता येते.

\begin{defn}\label{pt:seq:ex:defn:depth}
सूत्राची \emph{खोली}~$\depth{A}$ पुढीलप्रमाणे
परिभाषित करतो:
\begin{enumerate}
  \item $A$ अण्वीय असेल तर $\depth{A} = 0$.
  \item $A \ident \lnot B$, $A \ident \lforall[x][B]$,
  किंवा $A \ident \lexists[x][B]$ असेल तर
  $\depth{A} = \depth{B} + 1$.
  \item $A \ident (B \land C)$, $(B \lor C)$, किंवा $(B \lif
  C)$ असेल तर $\depth{A} = \max(\depth{B},\depth{C}) + 1$.
\end{enumerate}
\end{defn}

कोणत्याही पद~$t$ साठी $\depth{A(x)} =
\depth{A(t)}$ सिद्ध करणे कठीण नाही. आपल्याला
महत्त्वाचे हे की तार्किक नियमात सक्रिय विघटक सूत्रांपेक्षा प्रमुख
सूत्राची खोली नेहमी अधिक असते. संरचनात्मक
नियमांसाठी असा दावा नाही; तसेच संख्यापकाच्या काही
\Log{G3c} नियमांत जपलेली प्रमुख आस्थिती स्वतः त्याच
खोलीची असते. उदाहरणार्थ:
\begin{align*}
\depth{P(x)} = \depth{P(a)} & = 0,\\
\depth{\lforall[x][P(x)]} & = 1,\\
\depth{(\lforall[x][P(x)] \lif P(a))} & = \max(1,0) + 1 = 2.
\end{align*}
याचा वापर करून~$\depth{A}$ वर विगमनाने काही
निष्कर्ष सिद्ध करता येतात; पुढील निष्कर्ष त्यांपैकी एक.

\begin{prop}\label{pt:seq:ex:prop:G1c-axioms}
कोणत्याही~$A$ साठी, $A \Sequent A$ ची अशी
$\Log{G1c}$-सिद्धता आहे की तिच्यातील सर्व
स्वयंसिद्धके अण्वीय आहेत.
\end{prop}

\begin{proof}
  $\depth{A}$ वर विगमन करू. $\depth{A} = 0$
  असेल तर~$A$ अण्वीय आहे आणि $A \Sequent A$
  हीच अपेक्षित रूपाची \Log{G1c} मधील सिद्धता आहे.
  $\depth{A} > 0$ असेल तर~$A$ कमी खोलीच्या
  सूत्रांपासून बनते, उदा., $A \ident \lnot B$,
  $A \ident
  (B \land C)$, किंवा $A \ident \lforall[x][A(x)]$.
  विगमनात्मक गृहीतक असे: $\depth{D} < \depth{A}$
  असलेल्या प्रत्येक~$D$ साठी अण्वीय स्वयंसिद्धकांपासून
  $\Log{G1c} \Proves D \Sequent D$.

$A \ident \lnot B$ असेल, तर
  $\depth{B} < \depth{A}$; विगमनात्मक
  गृहीतकानुसार अण्वीय स्वयंसिद्धकांपासून
  $B \Sequent B$ ची \Log{G1c} मधील
  सिद्धता~$\pi_1$ आहे. तिचा पुढीलप्रमाणे
  $\lnot B \Sequent \lnot B$ च्या सिद्धता
  पर्यंत विस्तार करता येतो:
\[
  \AxiomC{}
  \RightLabel{$\pi_1$}
  \Deduce$B \fCenter B$
  \RightLabel{\LeftR{\lnot}}
  \UnaryInf$\lnot B, B \fCenter $
  \RightLabel{\RightR{\lnot}}
  \UnaryInf$\lnot B \fCenter \lnot B$
  \DisplayProof
\]

$A \ident (B \land C)$ असेल, तर
  $\depth{B} < \depth{A}$ आणि
  $\depth{C} < \depth{A}$. विगमनात्मक गृहीतकानुसार
  अण्वीय स्वयंसिद्धकांपासून $B
  \Sequent B$ आणि $C \Sequent C$ यांच्या
  \Log{G1c} मधील सिद्धता~$\pi_1$ आणि~$\pi_2$
  आहेत. त्यांचा पुढीलप्रमाणे
  $B \land C \Sequent B \land C$ च्या
  सिद्धता पर्यंत विस्तार करता येतो:
\[
  \AxiomC{}
  \RightLabel{$\pi_1$}
  \Deduce$B \fCenter B$
  \RightLabel{\LeftR{\Weakening}}
  \UnaryInf$B, C \fCenter B$
  \RightLabel{\LeftR{\land}}
  \UnaryInf$B \land C, C\fCenter B$
  \RightLabel{\LeftR{\land}}
  \UnaryInf$B \land C, B \land C \fCenter B$
  \AxiomC{}
  \RightLabel{$\pi_2$}
  \Deduce$C \fCenter C$
  \RightLabel{\LeftR{\Weakening}}
  \UnaryInf$B, C \fCenter C$
  \RightLabel{\LeftR{\land}}
  \UnaryInf$B \land C, C \fCenter C$
  \RightLabel{\LeftR{\land}}
  \UnaryInf$B \land C, B \land C \fCenter C$
  \RightLabel{\RightR{\land}}
  \BinaryInf$B \land C, B \land C \fCenter B \land C$
  \RightLabel{\LeftR{\Contraction}}
  \UnaryInf$B \land C \fCenter B \land C$
  \DisplayProof
  \]

आता $A \ident \lforall[x][B(x)]$ असे माना.
  $\lforall[x][B(x)]$ मध्ये न येणारा
  स्थिरांक~$c$ घ्या. मग $\depth{B(c)} =
  \depth{B(x)} < \depth{A}$; विगमनात्मक गृहीतकानुसार
  अण्वीय स्वयंसिद्धकांपासून $B(c)
  \Sequent B(c)$ ची \Log{G1c} मधील
  सिद्धता~$\pi_1$ आहे. तिचा पुढीलप्रमाणे
  $\lforall[x][B(x)] \Sequent \lforall[x][B(x)]$
  च्या सिद्धता पर्यंत विस्तार करता येतो:
\[
  \AxiomC{}
  \RightLabel{$\pi_1$}
  \Deduce$B(c) \fCenter B(c)$
  \RightLabel{\LeftR{\lforall}}
  \UnaryInf$\lforall[x][B(x)] \fCenter B(c)$
  \RightLabel{\RightR{\lforall}}
  \UnaryInf$\lforall[x][B(x)] \fCenter \lforall[x][B(x)]$
  \DisplayProof
\]

उरलेले प्रसंग प्रश्न म्हणून सोडले आहेत.
\end{proof}

\begin{prob}
  $A \ident (B \lor C)$, $A \ident (B \lif C)$
  आणि $A \ident \lexists[x][B(x)]$ या प्रसंगांचे
  काम करून \ref{pt:seq:ex:prop:G1c-axioms} ची सिद्धता पूर्ण करा.
\end{prob}

आपल्या कामासाठी $A \ident \lnot B$ या
प्रसंगात पुढील सिद्धता ही वापरता आली असती:
\[
\AxiomC{}
\RightLabel{$\pi_1$}
\Deduce$B \fCenter B$
\RightLabel{\RightR{\lnot}}
\UnaryInf$\fCenter B, \lnot B$
\RightLabel{\LeftR{\lnot}}
\UnaryInf$\lnot B \fCenter \lnot B$
\DisplayProof
\]
मात्र काही क्रमवर्ती पद्धतींमध्ये हे चालले नसते.
अंतःप्रज्ञावादी क्रमवर्ती कलन~\Log{G1i} हे
\Log{G1c} सारखेच आहे; फरक इतकाच की प्रत्येक
क्रमवर्तीच्या उजव्या बाजूस जास्तीत जास्त एकच
सूत्र ठेवता येते. $\Delta = \emptyset$
असेल, तर पहिली सिद्धता ही~\Log{G1i} मध्येही
सिद्धता असेल; परंतु दुसरी नसेल, कारण तिच्या
एका क्रमवर्तीच्या उजवीकडे $B$ आणि~$\lnot B$
दोन्ही आहेत.

$A \ident \lforall[x][B(x)]$ असताना
\Log{G1c} मध्येही नियमांचा दुसरा क्रम वापरता
आला नसता. पुढील रचना \emph{सिद्धता नाही}:
\[
  \AxiomC{}
  \RightLabel{$\pi_1$}
  \Deduce$B(c) \fCenter B(c)$
  \RightLabel{\RightR{\lforall}}
  \UnaryInf$B(c) \fCenter \lforall[x][B(x)]$
  \RightLabel{\LeftR{\lforall}}
  \UnaryInf$\lforall[x][B(x)] \fCenter \lforall[x][B(x)]$
  \DisplayProof
\]
कारण त्यात~$\RightR{\lforall}$ ची आयगेन चराची
अट मोडली जाते.
येथे~$x$ हे~$B(x)$ मध्ये खरोखर मुक्त आहे असे मानतो,
म्हणून~$c$ डावीकडील~$B(c)$ मध्ये येतो. संख्यापक
निष्फळ असेल आणि~$c$ त्या सूत्रात येतच नसेल, तर
दाखवलेल्या आयगेन-चर अटीचे उल्लंघन होत नाही.

\begin{prop}\label{pt:seq:ex:prop:G3c-axioms}
$A, \Gamma
\Sequent \Delta, A$ या रूपाची कोणतीही क्रमवर्ती
(~$A$ \emph{अण्वीय} असण्याची आवश्यकता नाही)
\Log{G3c} मध्ये सिद्धतायोग्य आहे.
\end{prop}

\begin{proof}
ही सिद्धता \ref{pt:seq:ex:prop:G1c-axioms} च्या सिद्धतेसारखीच
आहे, पण विगमनात्मक गृहीतकाबाबत काळजी घ्यावी
लागते. $n =
\depth{A} = 0$ हा प्रसंग पुन्हा सोपा आहे.
$n=0$ असेल, तर~$A$ अण्वीय आहे, आणि $A,
\Gamma \Sequent \Delta, A$ हे~\Log{G3c}
चे स्वयंसिद्धक आहे.

आता $n>0$ असे माना. पुन्हा प्रसंग वेगळे करू.
विगमनात्मक गृहीतक असे: $\depth{D} < n$ असलेल्या
सर्व~$D$ साठी आणि सर्व~$\Pi$ व~$\Lambda$
साठी $\Log{G3c} \Proves D, \Pi \Sequent \Lambda, D$.
$A \ident (B \land C)$ हा प्रसंग पाहू. विगमनात्मक
गृहीतक दोनदा वापरू: एकदा $D = B$,
$\Pi = C, \Gamma$, आणि $\Lambda
= \Delta$ घेऊन $B, C, \Gamma \Sequent
\Delta, B$ ची सिद्धता~$\pi_1$ मिळते; आणि
दुसऱ्यांदा $D = C$, $\Pi = B, \Gamma$, आणि
$\Lambda = \Delta$ घेऊन $B, C, \Gamma \Sequent
\Delta, C$ ची सिद्धता~$\pi_2$ मिळते.
$B \land C, \Gamma \Sequent
\Delta, B \land C$ ची सिद्धता पुढीलप्रमाणे:
\[
\AxiomC{}
\RightLabel{$\pi_1$}
\Deduce$B, C, \Gamma \fCenter \Delta, B$
\RightLabel{\LeftR{\land}}
\UnaryInf$B \land C, \Gamma \fCenter \Delta, B$
\AxiomC{}
\RightLabel{$\pi_2$}
\Deduce$B, C, \Gamma \fCenter \Delta, C$
\RightLabel{\LeftR{\land}}
\UnaryInf$B \land C, \Gamma \fCenter \Delta, C$
\RightLabel{\RightR{\land}}
\BinaryInf$B \land C, \Gamma \fCenter \Delta, B \land C$
\DisplayProof
\]

$A \ident \lforall[x][B(x)]$ साठी
~$\lforall[x][B(x)]$, $\Gamma$ किंवा~$\Delta$
मध्ये न येणारा स्थिरांक निवडा.
$D = B(c)$, $\Pi = \lforall[x][B(x)], \Gamma$
आणि $\Lambda = \Delta$ घेऊन विगमनात्मक गृहीतक
लावा. त्यातून $B(c),
\lforall[x][B(x)], \Gamma \Sequent \Delta, B(c)$
ची सिद्धता~$\pi_1$ मिळते; तिच्यापासून:
\[
\AxiomC{}
\RightLabel{$\pi_1$}
\Deduce$B(c), \lforall[x][B(x)], \Gamma \fCenter \Delta, B(c)$
\RightLabel{\LeftR{\lforall}}
\UnaryInf$\lforall[x][B(x)], \Gamma \fCenter \Delta, B(c)$
\RightLabel{\RightR{\lforall}}
\UnaryInf$\lforall[x][B(x)], \Gamma \fCenter \Delta, \lforall[x][B(x)]$
\DisplayProof
\]
आपण~$c$ ची केलेली निवड~$\RightR{\lforall}$
ची आयगेन चराची अट पूर्ण करते.
\end{proof}

\begin{prob}
  $A \ident (B \lif C)$ आणि $A \ident
  \lexists[x][B(x)]$ या प्रसंगांचे काम करून
  \ref{pt:seq:ex:prop:G3c-axioms} ची सिद्धता पुढे न्या.
\end{prob}













\section{नियमांचा अर्थ}\label{pt:seq:irl:sec}

क्रमवर्तीचा अर्थ आठवा: संरचना~$\Struct{M}$ मध्ये
क्रमवर्ती~$\Gamma \Sequent \Delta$ सत्य असते, जर
किमान एक $A \in \Gamma$ असत्य असेल किंवा किमान
एक~$A \in \Delta$ सत्य असेल. \Log{G3c} चे
तार्किक नियम त्यांच्या प्रमुख सूत्रांच्या
सत्यतेच्या अटी नियमबद्ध करतात असे वाचता येते.
\RightR{\lif} विचारात घ्या. त्याचे प्रमुख सूत्र
$A \lif B$ आहे. $A$ असत्य असेल किंवा~$B$
सत्य असेल तेव्हा ते सत्य असते. म्हणून क्रमवर्ती
$\ \Sequent A \lif B$ ही $A \Sequent B$
सत्य असेल तेव्हाच सत्य असते. या क्रमवर्तींच्या
पूर्वांगात आणि उत्तरांगात संदर्भातील सूत्रे
$\Gamma$, $\Delta$ जोडल्यास \RightR{\lif}
नियमाचे नेमके निष्कर्ष व पूर्वपक्ष मिळतात.
\LeftR{\lif} साठीही स्थिती अशीच आहे. $A
\lif B$ हे $A$~सत्य आणि $B$~असत्य असेल तेव्हा
असत्य असते. म्हणून $A \lif B
\Sequent \ $ ही क्रमवर्ती $\ \Sequent A$ आणि
$B \Sequent \ $ या दोन्ही सत्य असतील तेव्हाच
सत्य असते. \LeftR{\lif} चा निष्कर्ष त्याच्या
पूर्वपक्षांच्या संयुक्त सत्यतेशी समतुल्य आहे. या
रचनेमुळे \Log{G3c} चे नियम निर्दोष असतात (सर्व
पूर्वपक्ष सत्य असतील तर निष्कर्षही सत्य असतो).
त्यांच्यामुळे पूर्णताही मिळते.

खरे तर, कोणत्याही सत्यता-फलनीय संयोजकाच्या
सत्यतेच्या अटी अशा क्रमवर्तींच्या संचाने मांडता येतात.
कारण अभिजात तर्कशास्त्रातील प्रत्येक सत्यता-फलन
संयुक्त सामान्य रूपातील सूत्र ने दाखवता येते:
अण्वीय आणि निषेधित अण्वीय सूत्रे यांच्या
वियोजनांचे संयोजन. उदाहरणार्थ, $A
\oplus B$ हे $A$ आणि $B$ यांपैकी नेमके एकच
सत्य असेल तेव्हाच सत्य मानू; यालाच ``अनन्य किंवा''
म्हणतात. मग $A \oplus B$ हे $(A \lor B) \land
(\lnot A \lor \lnot B)$ सत्य असेल तेव्हाच सत्य
असते; आणि $(\lnot A \lor B)
\land (A \lor \lnot B)$ सत्य असेल तेव्हा ते असत्य
असते. म्हणून $\oplus$ साठी क्रमवर्ती-कलनाचे नियम
पुढीलप्रमाणे देता येतील:
\begin{align*}
&
\Axiom$\Gamma \fCenter \Delta, A, B$
\Axiom$A, B, \Gamma \fCenter \Delta$
\RightLabel{\RightR{\oplus}}
\BinaryInf$\Gamma \fCenter \Delta, A \oplus B$
\DisplayProof
\\
& \Axiom$A, \Gamma \fCenter \Delta, B$
\Axiom$B, \Gamma \fCenter \Delta, A$
\RightLabel{\LeftR{\oplus}}
\BinaryInf$A \oplus B, \Gamma \fCenter \Delta$
\DisplayProof
\end{align*}
हे नियमही निर्दोष व पूर्ण असतील. त्यांचे निष्कर्ष सर्व
पूर्वपक्ष सत्य असतील तेव्हाच सत्य असतात.

\begin{prob}
$A \downarrow B$ हे $A$ किंवा~$B$ यांपैकी
एकही सत्य नसल्यास आणि तरच सत्य आहे असे माना
(``NOR''). त्याचे \Log{G3c} मधील नियम कोणते असतील?
($A \downarrow B$ सत्य असेल तेव्हाच एकाच वेळी
सत्य असणाऱ्या दोन क्रमवर्ती शोधा. त्यातून
\RightR{\downarrow} नियम मिळेल. त्यानंतर $A
\downarrow B$ असत्य असेल तेव्हाच सत्य असणारी एक
क्रमवर्ती शोधा. त्यातून \LeftR{\downarrow} नियम मिळेल.)
\end{prob}

\Log{G3c} मध्ये प्रत्येक संयोजक आणि संख्यक
यांचे नेमके दोन नियम आहेत: एक डावीकडचा आणि एक
उजवीकडचा. येथे~$\lfalse$ हे स्वतंत्र तार्किक स्थिरांक
आहे; त्यासाठी डावीकडील स्वयंसिद्धक असते, उजवीकडील
स्वतंत्र नियम नसतो. पण \Log{G1c} मध्ये दोन
\LeftR{\land} आणि दोन \RightR{\lor} नियम आहेत.
याचे कारण असे की \Log{G1c} चे प्रतिमान असलेल्या
गेन्ट्झेनच्या मूळ~\Log{LK} पद्धतीला एक महत्त्वाच्या
अभिजातेतर तर्कशास्त्रासाठी, म्हणजे अंतःप्रज्ञावादी
तर्कशास्त्रासाठी,~\Log{LJ} हा प्रकार होता.
अंतःप्रज्ञावादी तर्कशास्त्रासाठी निर्दोष व पूर्ण
पद्धत मिळवण्याकरिता~\Log{LJ} मधील प्रत्येक
क्रमवर्तीच्या उत्तरांगात जास्तीत जास्त एकच
सूत्र ठेवण्याचे बंधन आहे. त्यामुळे
\Log{G3c} चा \RightR{\lor} नियम वापरता येत
नाही, कारण त्याच्या पूर्वपक्षाच्या उत्तरांगात $A$
आणि~$B$ ही दोन्ही येतात. म्हणून गेन्ट्झेनने
\Log{LJ} साठी \RightR{\lor} नियमांची जोडी
वापरली, आणि तीच~\Log{LK} साठीही वापरली.
तदनुरूप, अंतःप्रज्ञावादी प्रकार~\Log{G1i} हा
प्रत्येक उत्तरांगात जास्तीत जास्त एकच सूत्र
ठेवण्याचे बंधन घातलेली~\Log{G1c} पद्धत आहे.
(\Log{G3c} चाही अंतःप्रज्ञावादी प्रकार आहे;
मात्र त्यातील काही नियम \Log{G3c} मधील संबंधित
नियमांहून वेगळे आहेत. उदाहरणार्थ, त्यातही
\RightR{\lor} नियमांची जोडी वापरतात.)

\Log{G1c} चे संरचनात्मक नियम निर्दोष आहेत
हे सहज दिसते. उदाहरणार्थ, $\Gamma \Sequent \Delta$
च्या डावीकडे असत्य सूत्र किंवा उजवीकडे सत्य
सूत्र असेल, तर $\Gamma \Sequent
\Delta, A$ मध्येही तेच असते; $A$ सत्य आहे की
असत्य, याने फरक पडत नाही. म्हणून
\RightR{\Weakening} चा पूर्वपक्ष सत्य असेल तर
निष्कर्षही सत्य असतो. येथे उलट दिशेचा दावा मात्र
सत्य नाही: निष्कर्ष $A$ सत्य असल्यामुळे सत्य असू
शकतो; तेव्हा पूर्वपक्ष सत्य असेल याची खात्री नाही.

संरचनात्मक नियम मुळात का हवेत? उदाहरणार्थ,
$\ \Sequent A \lor \lnot A$ सिद्ध करण्यासाठी
आकुंचन (\RightR{\Contraction},
\LeftR{\Contraction}) लागते. $A \Sequent A$
पासून सुरुवात केल्यास, आधी~\RightR{\lnot} वापरून
$\Sequent A, \lnot A$ मिळते. मग
\RightR{\lor} च्या दोन आवृत्त्या क्रमाने वापरता येतात:
एकदा~$\lnot A$ पासून $A \lor \lnot
A$ मिळवण्यासाठी आणि एकदा~$A$ पासून
$A \lor \lnot A$ मिळवण्यासाठी. त्यातून $\ \Sequent A \lor \lnot
A, A \lor \lnot A$ मिळते. म्हणून \Log{G1c}
मध्ये $\Sequent A \lor \lnot A$ मिळवण्यासाठी
सिद्धता मधील त्याच सूत्राच्या दोन
उपस्थितींचे ``आकुंचन'' करता आले पाहिजे:
\[
\Axiom$A \fCenter A$
\RightLabel{\RightR{\lnot}}
\UnaryInf$\fCenter A, \lnot A$
\RightLabel{\RightR{\lor}}
\UnaryInf$\fCenter A, A \lor \lnot A$
\RightLabel{\RightR{\lor}}
\UnaryInf$\fCenter A \lor \lnot A, A \lor \lnot A$
\RightLabel{\RightR{\Contraction}}
\UnaryInf$\fCenter A \lor \lnot A$
\DisplayProof
\]
खरे तर, दुर्बलीकरण किंवा आकुंचन नसताना \Log{G1c}
मध्ये सिद्ध न करता येणाऱ्या वैध क्रमवर्ती आहेत.
उदाहरणार्थ,~$A$ आणि~$B$ ही वेगळी अणुरूप
सूत्रे घ्या. दुर्बलीकरण व आकुंचन वगळलेल्या तक्त्यातील
\Log{G1c} नियमांनी क्रमवर्ती
$A \Sequent B \lif A$ ची सिद्धता
\RightR{\lif} नेच संपली पाहिजे (कारण फक्त
$B \lif A$ प्रमुख असू शकते); त्यासाठी
$B, A \Sequent A$ हा पूर्वपक्ष हवा. तो
$A \Sequent A$ या स्वयंसिद्धकापासून मिळवायला
\LeftR{\Weakening} लागते:
\[
\Axiom$A \fCenter A$
\RightLabel{\LeftR{\Weakening}}
\UnaryInf$B, A \fCenter A$
\RightLabel{\RightR{\lif}}
\UnaryInf$A \fCenter B \lif A$
\DisplayProof
\]

\Log{G3c} मध्ये आकुंचन आणि दुर्बलीकरण
हे नियम अजिबात लागत नाहीत. दुर्बलीकरण आधीच
स्वयंसिद्धकांत अंतर्भूत आहे. दोन पूर्वपक्ष असलेल्या
नियमांमधील संदर्भ~$\Gamma$, $\Delta$ यांच्या दोन
प्रती एकत्र करून आकुंचन बहुतेक ठिकाणी टाळता येते.
\Log{G1c} आणि \Log{G3c} या दोन्हींमध्ये असे
``संदर्भ सामायिक करणारे'' नियम आहेत. एकच पूर्वपक्ष
असलेल्या नियमांत आकुंचन फक्त \RightR{\lor},
\LeftR{\land}, \RightR{\lexists} आणि~\LeftR{\lforall}
यांसाठी लागते. \Log{G3c} मधील या नियमांच्या
आवृत्त्या आकुंचन पूर्णपणे अनावश्यक करतात.
\Log{G2c} ही आणखी एक पद्धत आहे (
\ref{pt:seq:irl:tab:G2c} पाहा); तिच्या दोन-पूर्वपक्षी नियमांना
स्वतंत्र संदर्भ असतात. \Log{G2c} मध्ये आकुंचन
\Log{G1c} पेक्षाही अधिक महत्त्वाचे आहे.








\begin{table}
\begin{tabular}{@{}c@{\hskip-5ex}c@{}}
\multicolumn{2}{@{}c@{}}{स्वयंसिद्धके}\\[2ex]
$A \Sequent A$ & $\lfalse \Sequent \quad$
\\[2ex]
\multicolumn{2}{@{}c@{}}{संरचनात्मक नियम}\\[2ex]
\Axiom$ \Gamma \fCenter \Delta $
\RightLabel{\LeftR{\Weakening}}
\UnaryInf$ A, \Gamma \fCenter \Delta$
\DisplayProof
&
\Axiom$ \Gamma \fCenter \Delta$
\RightLabel{\RightR{\Weakening}}
\UnaryInf$ \Gamma \fCenter \Delta, A$
\DisplayProof
\\[4ex]
\Axiom$ A, A, \Gamma \fCenter \Delta $
\RightLabel{\LeftR{\Contraction}}
\UnaryInf$ A, \Gamma \fCenter \Delta$
\DisplayProof
&
\Axiom$ \Gamma \fCenter \Delta, A, A$
\RightLabel{\RightR{\Contraction}}
\UnaryInf$ \Gamma \fCenter \Delta, A$
\DisplayProof
\\[4ex]
\multicolumn{2}{@{}c@{}}{तार्किक नियम}\\[2ex]
\Axiom$ \Gamma \fCenter \Delta, A $
\RightLabel{\LeftR{\lnot}}
\UnaryInf$ \lnot A, \Gamma \fCenter \Delta$
\DisplayProof
&
\Axiom$A, \Gamma \fCenter \Delta$
\RightLabel{\RightR{\lnot}}
\UnaryInf$ \Gamma \fCenter \Delta, \lnot A $
\DisplayProof
\\[4ex]\begin{tabular}[t]{@{}l@{}}
\Axiom$ A, \Gamma \fCenter \Delta$
\RightLabel{\LeftR{\land}}
\UnaryInf$ A \land B, \Gamma \fCenter \Delta$
\DisplayProof
\\[3ex]
\Axiom$B, \Gamma \fCenter \Delta$
\RightLabel{\LeftR{\land}}
\UnaryInf$A \land B, \Gamma \fCenter \Delta$
\DisplayProof\\[3ex]
\end{tabular}
&
\Axiom$\Gamma_1 \fCenter \Delta_1, A$
\Axiom$ \Gamma_2 \fCenter \Delta_2, B$
\RightLabel{\RightR{\land}}
\BinaryInf$ \Gamma_1, \Gamma_2 \fCenter \Delta_1, \Delta_2, A \land B $
\DisplayProof
\\[4ex]\Axiom$A, \Gamma_1 \fCenter \Delta_1$
\Axiom$B, \Gamma_2 \fCenter \Delta_2$
\RightLabel{\LeftR{\lor}}
\BinaryInf$A \lor B, \Gamma_1, \Gamma_2 \fCenter \Delta_1, \Delta_2$
\DisplayProof
&
\begin{tabular}[t]{@{}r@{}}
\Axiom$\Gamma \fCenter \Delta, A$
\RightLabel{\RightR{\lor}}
\UnaryInf$ \Gamma \fCenter \Delta, A \lor B$
\DisplayProof
\\[3ex]
\Axiom$ \Gamma \fCenter \Delta, B$
\RightLabel{\RightR{\lor}}
\UnaryInf$ \Gamma \fCenter \Delta, A \lor B$
\DisplayProof\\[3ex]
\end{tabular}
\\[4ex]
\Axiom$ \Gamma_1 \fCenter \Delta_1, A$
\Axiom$ B, \Gamma_2 \fCenter \Delta_2$
\RightLabel{\LeftR{\lif}}
\BinaryInf$ A \lif B, \Gamma_1, \Gamma_2 \fCenter \Delta_1, \Delta_2$
\DisplayProof
&
\Axiom$ A, \Gamma \fCenter \Delta, B$
\RightLabel{\RightR{\lif}}
\UnaryInf$ \Gamma \fCenter \Delta, A \lif B $
\DisplayProof
\\[4ex]
\Axiom$ A(t), \Gamma \fCenter \Delta$
\RightLabel{\LeftR{\lforall}}
\UnaryInf$ \lforall[x][A(x)],\Gamma \fCenter \Delta$
\DisplayProof
&
\Axiom$ \Gamma \fCenter \Delta, A(a) $
\RightLabel{\RightR{\lforall}}
\UnaryInf$ \Gamma \fCenter \Delta, \lforall[x][A(x)]$
\DisplayProof
\\[4ex]
\Axiom$ A(a), \Gamma \fCenter \Delta $
\RightLabel{\LeftR{\lexists}}
\UnaryInf$ \lexists[x][A(x)], \Gamma \fCenter \Delta$
\DisplayProof
&
\Axiom$ \Gamma \fCenter \Delta, A(t) $
\RightLabel{\RightR{\lexists}}
\UnaryInf$ \Gamma \fCenter \Delta, \lexists[x][A(x)]$
\DisplayProof
\end{tabular}
\caption{\Log{G2c} चे नियम. $\RightR{\forall}$ आणि
$\LeftR{\exists}$ मध्ये स्थिरांक~$a$ हा
निष्कर्ष-क्रमवर्तीमध्ये येऊ नये. $\Gamma$ आणि
$\Delta$ हे सूत्रांचे बहुसंच आहेत.}
\label{pt:seq:irl:tab:G2c}
\end{table}
















\section{नियमित सिद्धता आणि प्रतिस्थापन}\label{pt:seq:qua:sec}

संख्यकांचे नियम काही गुंतागुंत निर्माण करतात;
कारण~\LeftR{\lexists} आणि~\RightR{\lforall} या
नियमांना ``आयगेन चराच्या अटी'' लागू होतात. या
बहुतेक अडचणी त्या निगमनांतील स्थिरांक~$a$
पुन्हा कधीही वापरला जाणार नाही याची खात्री करून
सोडवता येतात. म्हणून पुढील व्याख्या देऊ.

\begin{defn}
\Log{G1c} किंवा \Log{G3c} मधील सिद्धता~$\pi$
\emph{नियमित} असते, जर
\begin{enumerate}
  \item कोणताही आयगेन चर~$a$ एकाहून अधिक
  \LeftR{\lexists} किंवा~\RightR{\lforall} निगमनांचा
  आयगेन चर म्हणून वापरला जात नसेल; आणि
  \item $a$ हा \LeftR{\lexists} किंवा~\RightR{\lforall}
  निगमनाचा आयगेन चर असेल, तर तो~$\pi$ मध्ये
  फक्त त्या निगमनाच्या वरच येत असेल.
\end{enumerate}
\end{defn}

\begin{lem}\label{pt:seq:qua:lem:subst-regular} $\pi$ ही
नियमित असून तिची अंतिम क्रमवर्ती $\Gamma \Sequent
\Delta$ आहे असे माना. तसेच~$c$ हा~$\pi$ मध्ये
आयगेन चर म्हणून न वापरलेला स्थिरांक असो,
आणि~$t$ हे~$\pi$ चा कोणताही आयगेन चर नसलेले बंद पद असो.
येथील संख्यापक-नियमांत पद बंद असण्याची अट आहे;
त्या संकेतानुसार ही आदेशनाची मर्यादा आवश्यक आहे.
मग~$\pi$ मधील~$c$ ची प्रत्येक उपस्थिती~$t$ ने
बदलून मिळणारी~$\Subst{\pi}{t}{c}$ ही नियमित
सिद्धता असते. $\Subst{\pi}{t}{c}$ ची अंतिम क्रमवर्ती
$\Subst{\Gamma}{t}{c}
\Sequent \Subst{\Delta}{t}{c}$ असते; येथे
$\Subst{\Gamma}{t}{c} =
\Setabs{A(t)}{A(c) \in \Gamma}$.
\end{lem}

\begin{proof}
  $\pheight{\pi}$ वर विगमन करू. $\pheight{\pi} = 0$
  असेल तर ती फक्त एक स्वयंसिद्धक आहे. मग
  $\Subst{\pi}{t}{c}$ हेही स्वयंसिद्धक असते, कारण
  त्यातील कोणतेही सूत्र~$A(c)$ हे~$A(t)$
  होते.

$\pheight{\pi} > 0$ असेल तर~$\pi$
  मधील शेवटच्या निगमनानुसार प्रसंग वेगळे करू.
  संरचनात्मक नियमांचे (\Log{G1c} मध्ये) आणि
  विधानात्मक संयोजकांच्या तार्किक नियमांचे प्रसंग
  सोपे आहेत; ते प्रश्न म्हणून सोडतो. $\pi$ चा शेवट
  संख्यकाच्या निगमनाने होणारे प्रसंग पाहू.
  \begin{enumerate}
    \item शेवटचे निगमन $\RightR{\lforall}$ असेल, तर
    $\pi$ या रूपाची असते:
\[
    \AxiomC{}
    \RightLabel{$\pi'$}
    \Deduce$\Gamma(c) \fCenter \Delta'(c), A(a,c)$
    \RightLabel{\RightR{\lforall}}
    \UnaryInf$\Gamma(c) \fCenter \Delta'(c), \lforall[x][A(x,c)]$
    \DisplayProof
    \]
    विगमनात्मक गृहीतकानुसार $\Subst{\pi'}{t}{c}$
    ही $\Gamma(t) \fCenter \Delta'(t), A(a,t)$ ची
    नियमित सिद्धता आहे. $a$ हा~$\pi$ चा आयगेन
    चर असल्यामुळे~$a$ हे~$t$ मध्ये येत नाही.
    म्हणून $\Subst{\pi}{t}{c}$ मधील शेवटचे निगमन,
\[
    \AxiomC{}
    \RightLabel{$\Subst{\pi'}{t}{c}$}
    \Deduce$\Gamma(t) \fCenter \Delta'(t), A(a,t)$
    \RightLabel{\RightR{\lforall}}
    \UnaryInf$\Gamma(t) \fCenter \Delta'(t), \lforall[x][A(x,t)]$
    \DisplayProof
    \]
    योग्य आहे: त्याच्या निष्कर्षात~$a$ येत नाही.
    \item शेवटचे निगमन \Log{G3c} मधील
    $\LeftR{\lforall}$ असेल, तर $\pi$ चे रूप असे:
\[
    \AxiomC{}
    \RightLabel{$\pi'$}
    \Deduce$A(s(c),c), \lforall[x][A(x,c)], \Gamma(c) \fCenter \Delta'(c)$
    \RightLabel{\LeftR{\lforall}}
    \UnaryInf$\lforall[x][A(x,c)], \Gamma(c) \fCenter \Delta'(c)$
    \DisplayProof
    \]
    विगमनात्मक गृहीतकानुसार $\Subst{\pi'}{t}{c}$
    ही $A(s(t),t), \lforall[x][A(x,t)], \Gamma(t) \fCenter
    \Delta'(t)$ ची नियमित सिद्धता आहे.
    $\Subst{\pi}{t}{c}$ मधील शेवटचे निगमन,
\[
    \AxiomC{}
    \RightLabel{$\Subst{\pi'}{t}{c}$}
    \Deduce$A(s(t),t), \lforall[x][A(x,t)], \Gamma(t) \fCenter \Delta'(t)$
    \RightLabel{\LeftR{\lforall}}
    \UnaryInf$\lforall[x][A(x,t)], \Gamma(t) \fCenter \Delta'(t)$
    \DisplayProof
    \]
    योग्य आहे. \Log{G1c} मधील \LeftR{\lforall}
    साठीचा प्रसंगही असाच, पण काहीसा सोपा आहे.
    \item शेवटचे निगमन~\RightR{\lexists} किंवा
    \LeftR{\lexists} असेल, तर ते प्रश्न म्हणून सोडले आहे.
  \end{enumerate}
  प्रत्येक प्रसंगात आपण नियमित सिद्धता च्या
  शेवटी एक क्रमवर्ती जोडली (दोन पूर्वपक्ष असलेल्या
  नियमात दोन नियमित सिद्धता वापरल्या). नव्या
  अंतिम क्रमवर्तीमध्ये~$\pi$ च्या अंतिम
  क्रमवर्तीपेक्षा फक्त~$c$ च्या जागी पद~$t$ आलेले
  असते. तसेच $\pi$ आणि $\Subst{\pi}{t}{c}$
  यांचे आयगेन चर तेच राहतात; कारण~$t$ मध्ये
  $\pi$ चा कोणताही आयगेन चर नाही, नव्या अंतिम
  क्रमवर्तीमध्येही $\Subst{\pi}{t}{c}$ चा आयगेन चर
  येत नाही. त्यामुळे~$\Subst{\pi}{t}{c}$ नियमित असते.
\end{proof}

\begin{prop}
$\pi$ ही \Log{G1c} किंवा \Log{G3c} मधील
सिद्धता असेल, तर तिच्याशी तीच रचना
(विशेषतः तीच उंची) असलेली नियमित
सिद्धता~$\pi'$ मिळते, जी फक्त आयगेन चर
बदलल्यामुळे~$\pi$ पेक्षा वेगळी असते.
\end{prop}

\begin{proof}
फक्त या सिद्धतेपुरते,~$\pi$ मधील आयगेन चराचे
एखादे निगमन \emph{स्वच्छ} म्हणू, जर त्याचा आयगेन
चर इतर कोणत्याही निगमनाचा आयगेन चर नसेल आणि
त्या निगमनाच्या तत्काळ उप-सिद्धता व्यतिरिक्त
~$\pi$ मध्ये कुठेही येत नसेल; अन्यथा ते
\emph{अस्वच्छ} म्हणू. $\pi$ मधील अस्वच्छ आयगेन
चराच्या निगमनांची संख्या~$n$ असो. $n$ वर
विगमन करून~$\pi$ चे अपेक्षित नियमित
सिद्धता~$\pi'$ मध्ये रूपांतर करता येते हे
दाखवू. $n=0$ असेल तर~$\pi$ मधील आयगेन चराची
सर्व निगमने स्वच्छ आहेत; म्हणून~$\pi$ नियमित
असते आणि सिद्ध करण्यास काही उरत नाही.
अन्यथा, सर्वांत वरचे अस्वच्छ आयगेन-चर निगमन,
उदा.~\RightR{\lforall}, निवडा. त्यात संपणारी
~$\pi$ ची उप-सिद्धता~$\pi_1$ पुढील रूपाची आहे:
\[
    \AxiomC{}
    \RightLabel{$\pi_2$}
    \Deduce$\Gamma \fCenter \Delta, A(c)$
    \RightLabel{\RightR{\lforall}}
    \UnaryInf$\Gamma \fCenter \Delta, \lforall[x][A(x)]$
    \DisplayProof
    \]
$c$ हा आयगेन चर असल्यामुळे तो~$\Gamma$,
$\Delta$, किंवा $\lforall[x][A(x)]$ मध्ये
येत नाही. $\pi_2$ नियमित आहे, कारण तिच्यातील
कोणतेही आयगेन-चर निगमन अस्वच्छ नाही.
~$\pi$ मध्ये न येणारा स्थिरांक~$a$ घ्या.
\ref{pt:seq:qua:lem:subst-regular} नुसार
$\Subst{\pi_2}{a}{c}$ ही $\Gamma \fCenter \Delta,
A(a)$ ची नियमित सिद्धता आहे; तिच्यावर
~$\RightR{\lforall}$ लावता येतो. $\pi$ मधील
~$\pi_1$ च्या जागी पुढील सिद्धता~$\pi_1'$ ठेवू:
\[
    \AxiomC{}
    \RightLabel{$\Subst{\pi_2}{a}{c}$}
    \Deduce$\Gamma \fCenter \Delta, A(a)$
    \RightLabel{\RightR{\lforall}}
    \UnaryInf$\Gamma \fCenter \Delta, \lforall[x][A(x)]$
    \DisplayProof
    \]
त्यामुळे एक अस्वच्छ आयगेन-चर निगमन स्वच्छ बनते;
फरक फक्त आयगेन चर~$c$ च्या जागी~$a$ आल्याचा
आहे. परिणामी सिद्धतेतील अस्वच्छ आयगेन-चर
निगमनांची संख्या~$n-1$ पेक्षा जास्त राहत नाही,
म्हणून विगमनात्मक गृहीतकातून निष्कर्ष मिळतो.
\end{proof}

नुकताच सिद्ध केलेला निष्कर्ष आपण स्पष्टपणे
संदर्भ म्हणून वापरणार नाही; पण पुढील सर्व
सिद्धता नियमित आहेत असे मानू. त्या नियमित
नसल्या तर आयगेन चर बदलून त्यांना नेहमी नियमित
बनवता येते.













\section{अनुज्ञेय आणि व्युत्पाद्य नियम}\label{pt:seq:adm:sec}

सिद्धता-उपपत्तीतील अनेक निष्कर्ष एखादा विशिष्ट नियम वापरून
सिद्ध करता येणारी प्रत्येक गोष्ट तो नियम न वापरताही
सिद्ध करता येते का, यासंबंधी असतात किंवा अशा
निष्कर्षांचा वापर करतात. कट-निरसन प्रमेय हे याचे सर्वांत
प्रसिद्ध उदाहरण आहे. त्यानुसार,~\Cut{} असलेल्या क्रमवर्ती
पद्धतीचे नियम वापरून सिद्ध करता येणारी कोणतीही क्रमवर्ती
~\Cut न वापरता सिद्ध करता येते. हा गुणधर्म इतका
महत्त्वाचा आहे की त्याला स्वतंत्र नाव आहे.

\begin{defn}
एखाद्या क्रमवर्ती पद्धतीत नियम~$R$ \emph{अनुज्ञेय} असतो,
जर त्या पद्धतीत नियम~$R$ जोडल्यावर सिद्ध करता येणारी
प्रत्येक क्रमवर्ती~$R$ शिवायही सिद्ध करता येत असेल.
\end{defn}

\begin{ex}\label{pt:seq:adm:ex:land-deriv}
नियम~$\LeftR{\land}$ चे \Log{G1c} आणि \Log{G3c} मध्ये
वेगवेगळे रूप आहे. \Log{G1c} मध्ये या नियमाच्या दोन आवृत्त्या
आहेत: एकीत प्रमुख सूत्र~$A \land B$ मधील
डावा संयोज्य~$A$ पूर्वपक्षात असतो, तर दुसरीत उजवा
संयोज्य~$B$ असतो:
\[
\Axiom$ A, \Gamma \fCenter \Delta$
\RightLabel{$\LeftR{\land}_1$}
\UnaryInf$ A \land B, \Gamma \fCenter \Delta$
\DisplayProof
\qquad
\Axiom$B, \Gamma \fCenter \Delta$
\RightLabel{$\LeftR{\land}_2$}
\UnaryInf$A \land B, \Gamma \fCenter \Delta$
\DisplayProof
\]
याउलट, \Log{G3c} मध्ये एकच नियम आहे:
\[
\Axiom$ A, B, \Gamma \fCenter \Delta$
\RightLabel{$\LeftR{\land}_3$}
\UnaryInf$ A \land B, \Gamma \fCenter \Delta$
\DisplayProof
\]
आता ``\Log{G3c} चा नियम~$\LeftR{\land}_3$
हा \Log{G1c} मध्ये अनुज्ञेय आहे का?'' असा प्रश्न विचारता येईल.
उत्तर होय आहे:~$\LeftR{\land}_3$ वापरून केलेल्या कोणत्याही
निगमनाचे \Log{G1c} च्याच नियमांनी थेट अनुकरण करता येते.
त्यासाठी~$\LeftR{\land}_1$ आणि $\LeftR{\land}_2$
यांबरोबर संरचनात्मक नियम~$\LeftR{\Contraction}$ ही लागेल:
\[
\Axiom$A, B, \Gamma \fCenter \Delta$
\RightLabel{$\LeftR{\land}_1$}
\UnaryInf$A \land B, B, \Gamma \fCenter \Delta$
\RightLabel{$\LeftR{\land}_2$}
\UnaryInf$A \land B, A \land B, \Gamma \fCenter \Delta$
\RightLabel{\LeftR{\Contraction}}
\UnaryInf$A \land B, \Gamma \fCenter \Delta$
\DisplayProof
\]
\end{ex}

एका नियमाने केलेल्या निगमनाचे इतर नियमांनी
अनुकरण करणे हा नियम अनुज्ञेय असल्याचे दाखवण्याचा एक मार्ग
आहे. पण तो एकमेव मार्ग नाही; काही नियम अनुज्ञेय असूनही
अशा प्रकारे त्यांचे अनुकरण करता येत नाही. ज्या नियमांचे
असे अनुकरण करता येते त्यांना \emph{व्युत्पाद्य} म्हणतात.

\begin{defn}
एखादा नियम~$R$,
\[
\AxiomC{$S_1$}
\AxiomC{$\dots$}
\AxiomC{$S_n$}
\RightLabel{$R$}
\TrinaryInfC{$S$}
\DisplayProof
\]
पद्धतीचे नियम आणि स्वयंसिद्धके, तसेच $S_1$, \dots,~$S_n$
या रूपांच्या अतिरिक्त प्रारंभीच्या क्रमवर्ती वापरून
क्रमवर्ती~$S$ ची रूपरेषात्मक सिद्धता असेल, तर तो नियम
त्या पद्धतीत \emph{व्युत्पाद्य} म्हणतात.
\end{defn}

\ref{pt:seq:adm:ex:land-deriv} मधील सिद्धता दाखवते
की~$\LeftR{\land}_3$ हा नियम \Log{G1c} मध्ये व्युत्पाद्य
आहे. कारण ती~$\LeftR{\land}_3$ च्या पूर्वपक्षांना प्रारंभीच्या
क्रमवर्ती मानून, फक्त \Log{G1c} चे नियम वापरून,
$\LeftR{\land}_3$ च्या निष्कर्षाची रूपरेषात्मक सिद्धता
आहे. (त्या सिद्धता मध्ये \Log{G1c}
मधील कोणतेही स्वयंसिद्धक वापरलेले नाही.)

\begin{prop}\label{pt:seq:adm:prop:lor-G3-adm}
\Log{G3c} चे नियम~$\LeftR{\land}$ आणि $\RightR{\lor}$
हे \Log{G1c} मध्ये व्युत्पाद्य आहेत.
\end{prop}

\begin{proof}
  $\LeftR{\land}$ साठीची सिद्धता \ref{pt:seq:adm:ex:land-deriv} मध्ये
  दिली आहे. $\RightR{\lor}$ साठीची सिद्धता प्रश्न म्हणून सोडली आहे.
\end{proof}

\begin{prob}
\Log{G3c} चा नियम~$\RightR{\lor}$ हा \Log{G1c} मध्ये
व्युत्पाद्य आहे हे दाखवा.
\end{prob}

\begin{prob}
$\liff$ हा परिभाषित संकारक आहे असे माना:
$A \liff B$ हे $(A \lif B) \land (B \lif A)$
याचे संक्षिप्त रूप आहे. पुढील नियम
\begin{enumerate}
\item \Axiom$\Gamma, A, B \fCenter \Delta$
\Axiom$\Gamma \fCenter \Delta, A, B$
\RightLabel{\LeftR{\liff}}
\BinaryInf$A \liff B, \Gamma \fCenter \Delta$
\DisplayProof
\item
\Axiom$A, \Gamma \fCenter \Delta, B$
\Axiom$B, \Gamma \fCenter \Delta, A$
\RightLabel{\RightR{\liff}}
\BinaryInf$\Gamma \fCenter \Delta, A \liff B$
\DisplayProof
\end{enumerate}
(a)~\Log{G1c} आणि (b)~\Log{G3c} मध्ये व्युत्पाद्य
आहेत हे दाखवा.
\end{prob}

\begin{ex}\label{pt:seq:adm:ex:weak-adm}
\Log{G1c} चा नियम~$\RightR{\Weakening}$,
\[
\Axiom$ \Gamma \fCenter \Delta$
\RightLabel{\RightR{\Weakening}}
\UnaryInf$ \Gamma \fCenter \Delta, A$
\DisplayProof
\]
हा \Log{G3c} मध्ये अनुज्ञेय आहे. सिद्धतेची कल्पना सोपी
आहे: \Log{G3c} मध्ये पूर्वपक्ष~$\Gamma \Sequent \Delta$
याची सिद्धता~$\pi$ आहे असे माना. $\pi$ मधील
प्रत्येक क्रमवर्तीच्या निष्कर्षभागात सूत्र~$A$ जोडा.
स्वयंसिद्धके स्वयंसिद्धकेच राहतात; योग्य निगमनेही योग्यच
राहतात. नव्या सिद्धता मध्ये हे सूत्र सर्वत्र
जोडलेले असते; तिची अंतिम क्रमवर्ती
$\Gamma \Sequent \Delta, A$ असते.
यापूर्वी~$\pi$ चे आयगेन स्थिरांक~$A$ मध्ये येणाऱ्या
स्थिरांकांपासून वेगळे नवी नावे देऊन घ्यायचे. नावबदल
अंतिम क्रमवर्ती आणि उंची जपतो; त्यानंतरच~$A$ सर्व
क्रमवर्तींमध्ये जोडायचे, जेणेकरून संख्यापक-नियमांची
आयगेन-चर अट जपली जाते.

तथापि, नियम~$\RightR{\Weakening}$ हा
\Log{G3c} मध्ये व्युत्पाद्य नाही. समजा तो व्युत्पाद्य
आहे; म्हणजे \Log{G3c} ची स्वयंसिद्धके व नियम आणि
कदाचित अतिरिक्त प्रारंभीची क्रमवर्ती~$\Gamma \Sequent
\Delta$ वापरून $\Gamma \Sequent \Delta, A$ ची
सिद्धता मिळते. $\RightR{\Weakening}$ व्युत्पाद्य
असण्यासाठी आवश्यक असल्याप्रमाणे ही सिद्धता पूर्णतः
रूपरेषात्मक असेल, तर $\Gamma$ आणि~$\Delta$ काढून
टाकून तिच्यापासून~$\Sequent A$ ची सिद्धता
बनवता येईल. त्यामुळे अतिरिक्त प्रारंभीच्या क्रमवर्ती
$\Gamma \Sequent \Delta$ चे रूपांतर रिकाम्या
क्रमवर्ती~$\quad \Sequent \quad$ मध्ये होईल. रिकाम्या
क्रमवर्तीमध्ये सूत्रे अजिबात नसतात, तर \Log{G3c}
च्या प्रत्येक नियमाच्या पूर्वपक्षात किमान एक सक्रिय
सूत्र आवश्यक असते. म्हणून या रूपरेषात्मक सिद्धतेतील
कोणत्याही निगमनाचा पूर्वपक्ष रिकामी क्रमवर्ती असू शकत
नाही. परिणामी, अतिरिक्त प्रारंभीची क्रमवर्ती~$\Gamma
\Sequent \Delta$ प्रत्यक्ष न वापरल्यासच ती सिद्धता
रूपरेषात्मक असू शकते. दुसऱ्या शब्दांत,
$\Gamma \Sequent \Delta, A$ चे \Log{G3c} चीच
स्वयंसिद्धके व नियम वापरून रूपरेषात्मक सिद्धता मिळेल;
मग या रूपाची कोणतीही क्रमवर्ती \Log{G3c} मध्ये
सिद्धता असलेली ठरेल. पण हे अशक्य आहे: \Log{G3c}
निर्दोष असल्यामुळे प्रत्येक संरचनांमध्ये सत्य
असलेल्या क्रमवर्तीच ते सिद्ध करू शकते. उदाहरणार्थ,
$\Gamma = \Delta = \emptyset$ आणि $A \ident \lfalse$
घेतल्यास, \Log{G3c} ला $\quad \Sequent \lfalse$
ही क्रमवर्ती सिद्ध करावी लागेल; ते तसे करत नाही.
\end{ex}

आता \Log{G3c} मध्ये दुर्बलीकरण हा अनुज्ञेय
नियम आहे, याची नीट विगमनात्मक सिद्धता देऊ. खरे तर,
वरील कल्पनेप्रमाणे किंचित अधिक प्रबळ निष्कर्ष मिळतो:
$\pi$ मधील प्रत्येक क्रमवर्तीच्या उजवीकडे~$A$
जोडल्याने सिद्धता ची रचना बदलत नाही; विशेषतः तिची
उंची वाढत नाही. हे महत्त्वाचे असल्याने स्वतंत्र व्याख्या
देऊ.

\begin{defn}
एखादा नियम~$R$,
\[
\AxiomC{$S_1$}
\AxiomC{$\dots$}
\AxiomC{$S_n$}
\RightLabel{$R$}
\TrinaryInfC{$S$}
\DisplayProof
\]
पद्धतीत \emph{उंची जपून अनुज्ञेय} (किंवा
\emph{उंची जपून-अनुज्ञेय}) असतो, जर $i = 1$, \dots,~$n$
साठी $\Proves[h_i] S_i$ असताना $m = \max(h_1, \dots,
h_n)$ घेऊन $\Proves_m S$ मिळत असेल.
\end{defn}

नियम उंची जपून अनुज्ञेय असण्यासाठी पुढील अट
पुरेशी आहे: पूर्वपक्ष~$S_i$ यांच्या उंची~$h_i =
\pheight{\pi_i}$ असलेल्या सिद्धता~$\pi_i$ दिल्या
असता, निष्कर्षाची अशी सिद्धता~$\pi$ असावी की
तिची उंची~$\pheight{\pi}$ ही~$h_i$ यांच्या कमाल
मूल्यापेक्षा मोठी नाही. विशेषतः, एकच पूर्वपक्ष~$S_1$
असेल तर $\pheight{\pi} \le \pheight{\pi_1}$ पुरेसे
आहे. $\RightR{\Weakening}$ आणि $\LeftR{\Weakening}$
बाबतीत तर प्रत्यक्षात $\pheight{\pi} =
\pheight{\pi_1}$ असते; पण समानता आवश्यक नाही.

\begin{prop}\label{pt:seq:adm:prop:weak-G3c-adm}
\Log{G1c} चे नियम~$\RightR{\Weakening}$ आणि
$\LeftR{\Weakening}$ हे \Log{G3c} मध्ये
उंची जपून अनुज्ञेय आहेत.
\end{prop}

\begin{proof}
  $\pheight{\pi} = n$ उंची असलेली सिद्धता~$\pi$
  वापरून $\Log{G3c} \Proves \Gamma \Sequent \Delta$
  असेल, तर $\pheight{\pi'} = n$ असलेली
  $\Gamma \Sequent \Delta, A$ ची \Log{G3c}-सिद्धता
  $\pi'$ आहे, हे दाखवायचे आहे. $n$ वर विगमन करू.
  $n=0$ असेल तर~$\pi$ हे फक्त $\Gamma \Sequent
  \Delta$ स्वयंसिद्धक आहे (म्हणजे एखादे अण्वीय~$C$
  हे~$\Gamma$ आणि~$\Delta$ दोन्हीत येते, किंवा
  $\lfalse$ डावीकडे असते); म्हणून
  $\Gamma \Sequent \Delta, A$ हेही स्वयंसिद्धक आहे.
  $n = \pheight{\pi} >0$ असेल तर सिद्धता~$\pi$
  मध्ये शेवटचे एक निगमन असते. कोणता नियम वापरला यानुसार
  उदाहरणे वेगळी करतो. एकच उदाहरण देऊ: शेवटचे निगमन
  $\RightR{\land}$ होते असे समजा. मग $\Delta =
  \Delta', B \land C$; शेवटच्या निगमनातील तत्काळ
  उप-सिद्धता~$\pi_1$ आणि~$\pi_2$ यांच्या अंतिम
  क्रमवर्ती अनुक्रमे $\Gamma \Sequent \Delta', B$
  आणि $\Gamma \Sequent \Delta', C$ असतात. म्हणजे
  $\pi$ चे रूप असे असते:
\[
  \AxiomC{}
  \RightLabel{$\pi_1$}
  \Deduce$\Gamma \fCenter \Delta', B$
  \AxiomC{}
  \RightLabel{$\pi_2$}
  \Deduce$\Gamma \fCenter \Delta', C$
  \RightLabel{\RightR{\land}}
  \BinaryInf$\Gamma \fCenter \Delta', B \land C$
  \DisplayProof
  \]
  तसेच $n =
  \max(\pheight{\pi_1}, \pheight{\pi_2}) + 1$.
  म्हणून $\pheight{\pi_1} <
  n$ आणि $\pheight{\pi_2} < n$; त्यामुळे विगमनात्मक
  गृहीतक लागू होते: $\Gamma \Sequent
  \Delta', B, A$ आणि $\Gamma \Sequent \Delta', C, A$
  यांच्या अनुक्रमे सिद्धता~$\pi_1'$ आणि~$\pi_2'$
  मिळतात. $\RightR{\land}$ नियमाने $\Gamma
  \Sequent \Delta', B \land C, A$ ची सिद्धता~$\pi'$
  मिळते:
\[
  \AxiomC{}
  \RightLabel{$\pi_1'$}
  \Deduce$\Gamma \fCenter \Delta', B, A$
  \AxiomC{}
  \RightLabel{$\pi_2'$}
  \Deduce$\Gamma \fCenter \Delta', C, A$
  \RightLabel{\RightR{\land}}
  \BinaryInf$\Gamma \fCenter \Delta', B \land C, A$
  \DisplayProof
  \]
  आपल्याला $\pheight{\pi'} =
  \max(\pheight{\pi_1'}, \pheight{\pi_2'}) + 1 = \max(\pheight{\pi_1},
  \pheight{\pi_2}) + 1 = \pheight{\pi}$ मिळते.
\end{proof}

\ref{pt:seq:adm:prop:weak-G3c-adm} अत्यंत उपयोगी आहे.
त्याचा वारंवार वापर दाखवण्यासाठी एक संकेतन ठरवू:
$\pi$ ही~$\Gamma \Sequent \Delta$ ची सिद्धता
असेल, तर $\pi[\Pi \Sequent \Lambda]$ म्हणजे डावीकडे
$\Pi$ मधील सर्व सूत्रे आणि उजवीकडे~$\Lambda$
मधील सर्व सूत्रे जोडण्यासाठी दुर्बलीकरणाच्या
अनुज्ञेयतेचा केलेला वापर. परिणामी ती $\Gamma, \Pi
\Sequent \Delta, \Lambda$ ची सिद्धता असते.













\section{नियमांची व्युत्क्रमणीयता}\label{pt:seq:inv:sec}

\begin{defn}
एखादा नियम~$R$,
\[
\AxiomC{$S_1$}
\AxiomC{$\dots$}
\AxiomC{$S_n$}
\RightLabel{$R$}
\TrinaryInfC{$S$}
\DisplayProof
\]
पद्धतीत \emph{व्युत्क्रमणीय} असतो, जर $\Proves S$
असताना $i = 1$, \dots,~$n$ साठी $\Proves[h_i] S_i$
मिळत असेल. तो \emph{उंची जपून व्युत्क्रमणीय}
(किंवा \emph{उंची जपून-व्युत्क्रमणीय}) असतो, जर
$\Proves_h S$ असताना $i = 1$, \dots,~$n$
साठी $\Proves[h] S_i$ मिळत असेल.
\end{defn}

\begin{lem}\label{pt:seq:inv:lem:G3c-invert}
\Log{G3c} मधील $\lnot$, $\land$, $\lor$, आणि $\lif$
यांचे नियम उंची जपून व्युत्क्रमणीय आहेत;
म्हणजे:
\begin{enumerate}
  \item \RightR{\lnot}: $\Log{G3c} \Proves[n] \Gamma \Sequent
  \Delta, \lnot A$ असेल, तर $\Log{G3c} \Proves[n] A, \Gamma \Sequent
  \Delta$.
  \item \LeftR{\lnot}: $\Log{G3c} \Proves[n] \Gamma, \lnot A
  \Sequent \Delta$ असेल, तर $\Log{G3c} \Proves[n] \Gamma \Sequent \Delta,
  A$.
  \item \RightR{\land}: $\Log{G3c} \Proves[n] \Gamma \Sequent
  \Delta, A \land B$ असेल, तर $\Log{G3c} \Proves[n] \Gamma \Sequent
  \Delta, A$ आणि $\Log{G3c} \Proves[n] \Gamma \Sequent
  \Delta, B$.
  \item \LeftR{\land}: $\Log{G3c} \Proves[n] \Gamma, A \land B
  \Sequent \Delta$ असेल, तर $\Log{G3c} \Proves[n] A, B, \Gamma \Sequent
  \Delta$.
  \item \RightR{\lor}: $\Log{G3c} \Proves[n] \Gamma \Sequent
  \Delta, A \lor B$ असेल, तर $\Log{G3c} \Proves[n] \Gamma \Sequent
  \Delta, A, B$.
  \item \LeftR{\lor}: $\Log{G3c} \Proves[n] A \lor B, \Gamma \Sequent
  \Delta$ असेल, तर $\Log{G3c} \Proves[n] A, \Gamma \Sequent
  \Delta$ आणि $\Log{G3c} \Proves[n] B, \Gamma \Sequent
  \Delta$.
  \item \RightR{\lif}: $\Log{G3c} \Proves[n] \Gamma \Sequent
  \Delta, A \lif B$ असेल, तर $\Log{G3c} \Proves[n] A, \Gamma \Sequent
  \Delta, B$.
  \item \LeftR{\lif}: $\Log{G3c} \Proves[n] A \lif B, \Gamma
  \Sequent \Delta$ असेल, तर $\Log{G3c} \Proves[n] \Gamma \Sequent \Delta,
  A$ आणि $\Log{G3c} \Proves[n] B, \Gamma \Sequent \Delta$.
\end{enumerate}
\end{lem}

\begin{proof}
वरील प्रमेयिकेत दिलेल्या विधानीय नियम~$R$ साठी पुढील
गोष्ट दाखवायची आहे:~$R$ च्या निष्कर्ष~$S$ ची
सिद्धता~$\pi$ असेल, तर $R$ च्या पूर्वपक्ष~$S_i$
यांच्या सिद्धता~$\pi_i$ सुद्धा असतात, आणि
$\pheight{\pi_i}
\le \pheight{\pi}$. नियम~\LeftR{\land} विचारात घ्या:
अंतिम क्रमवर्ती $A \land B, \Gamma
\Sequent \Delta$ या रूपाची असलेली सिद्धता~$\pi$
आहे असे समजा. $A, B, \Gamma \Sequent \Delta$
ची अशी सिद्धता~$\pi'$ आहे की $\pheight{\pi'} \le
\pheight{\pi}$, हे दाखवायचे आहे. $n = \pheight{\pi}$
वर विगमन करून हे सिद्ध करू.

$n = 0$ असेल तर $\pi$ फक्त एकच स्वयंसिद्धक
असतो: $A \land B, \Gamma
\Sequent \Delta$. \Log{G3c} मधील स्वयंसिद्धके
$C, \Gamma
\Sequent \Delta, C$ या रूपाची असतात, आणि~$C$
अण्वीय असतो. म्हणजेच $C$ हा $A \land B$
असू शकत नाही. $A \land B, \Gamma \Sequent \Delta$
स्वयंसिद्धक असेल, तर एखादा अण्वीय~$C$ हा~$\Gamma$
आणि~$\Delta$ दोहोंचा घटक असला पाहिजे.
मग $A, B, \Gamma \Sequent \Delta$ हाही
स्वयंसिद्धक असतो आणि आपल्याला हवी असलेली
सिद्धता~$\pi'$ मिळाली (तिची उंची ही~$0$).
दुसरा स्वयंसिद्धक-प्रसंग म्हणजे~$\lfalse\in\Gamma$.
संयोगाची आस्थिती बदलल्यावरही तो डावीकडे राहतो;
म्हणून अपेक्षित क्रमवर्ती पुन्हा उंची~$0$ चे
असत्य-स्वयंसिद्धक आहे.

$n>0$ असेल, तर सिद्धतेचा शेवट निगमनाने
होतो. $n$ साठी दावा सिद्ध करायचा आहे; परंतु संदर्भ
वेगळा असला तरी उंची~$<n$ असलेल्या कोणत्याही
सिद्धता साठी दावा खरा आहे असे मानता येते. म्हणजे
कोणत्याही $k<n$ आणि कोणत्याही $\Pi$ व $\Lambda$
साठी $A \land B, \Pi \Sequent \Lambda$ ची
सिद्धता उंची~$k$ असेल, तर $A, B, \Pi
\Sequent \Lambda$ ची सिद्धता उंची~$\le k$
असलेली मिळते.

दोन प्रसंग आहेत: डावीकडील $A \land
B$ प्रमुख असतो किंवा नसतो. तो प्रमुख असेल तर
शेवटचे निगमन~\LeftR{\land} असते आणि तत्काळ
उप-सिद्धता~$\pi_1$ ची अंतिम क्रमवर्ती
$A, B, \Gamma \Sequent \Delta$ असते. या
प्रसंगात निष्कर्ष मिळतो, कारण $\pheight{\pi_1} <
\pheight{\pi}$.

$A \land B$ प्रमुख नसेल तर दुसरे
एखादे सूत्र प्रमुख असते आणि $A
\land B$ हा शेवटच्या निगमनाच्या संदर्भात असतो.
मग शेवटच्या नियम~$R$ च्या तत्काळ उप-सिद्धता
ना विगमनात्मक गृहीतक लावता येते. त्यातून~$\pi$
च्या तत्काळ उप-सिद्धता च्या उंचींपेक्षा मोठी
उंची नसलेली सिद्धता~$\pi_1'$ किंवा
दोन सिद्धता~$\pi_1'$ आणि~$\pi_2'$ मिळतात.
$\pi_1'$ आणि~$\pi_2'$ यांच्या अंतिम क्रमवर्तींमध्ये
डावीकडील संदर्भात $A \land B$ ऐवजी $A, B$
येतात. त्यांना नियम~$R$ लावून सिद्धता~$\pi'$
मिळते. उदाहरणार्थ,~$\pi$ चा शेवट~\LeftR{\lif}
ने होत असेल तर:
\[
\AxiomC{}
\RightLabel{$\pi_1$}
\Deduce$A \land B, \Gamma' \fCenter \Delta, C$
\AxiomC{}
\RightLabel{$\pi_2$}
\Deduce$A \land B, \Gamma', D \fCenter \Delta$
\RightLabel{\LeftR{\lif}}
\BinaryInf$A \land B, \Gamma', C \lif D \fCenter \Delta$
\DisplayProof
\]
येथे डावीकडील $C \lif D$ प्रमुख असून डावीकडील
संदर्भ $A \land B, \Gamma'$ आहे (म्हणजे
$\Gamma = \Gamma', C \lif D$). विगमनात्मक
गृहीतकाने अनुक्रमे $A, B, \Gamma' \Sequent \Delta, C$
आणि $A, B, \Gamma', D \Sequent \Delta$
यांच्या सिद्धता~$\pi_1'$
आणि~$\pi_2'$ मिळतात. या दोन क्रमवर्ती पुन्हा
\LeftR{\lif} चे पूर्वपक्ष बनून~$\pi'$ देतात:
\[
\AxiomC{}
\RightLabel{$\pi_1'$}
\Deduce$A, B, \Gamma' \fCenter \Delta, C$
\AxiomC{}
\RightLabel{$\pi_2'$}
\Deduce$A, B, \Gamma', D \fCenter \Delta$
\RightLabel{\LeftR{\lif}}
\BinaryInf$A, B, \Gamma', C \lif D \fCenter \Delta$
\DisplayProof
\]
आपल्याला पुढील मिळते:
\begin{align*}
\pheight{\pi} & = \max(\pheight{\pi_1}, \pheight{\pi_2}) + 1\\
\pheight{\pi'} & = \max(\pheight{\pi_1'}, \pheight{\pi_2'}) + 1
\intertext{व्याख्येनुसार; आणि}
\pheight{\pi_1'} & \le \pheight{\pi_1}\\
\pheight{\pi_2'} & \le \pheight{\pi_2}
\intertext{विगमन गृहीतकानुसार. म्हणून,}
\pheight{\pi'} & \le \pheight{\pi}.
\end{align*}
\end{proof}

\begin{prob}
\RightR{\lif} साठी \ref{pt:seq:inv:lem:G3c-invert} ची सिद्धता द्या.
\end{prob}

\begin{lem}\label{pt:seq:inv:lem:invert-quant}
\RightR{\lforall} आणि \LeftR{\lexists} हे नियम
पुढील अर्थाने उंची जपून व्युत्क्रमणीय आहेत:
\begin{enumerate}
  \item $\Log{G3c} \Proves[n] \Gamma \Sequent \Delta,
  \lforall[x][A(x)]$ असेल, तर $\Log{G3c} \Proves[n] \Gamma \Sequent
  \Delta, A(c)$; आणि
  \item $\Log{G3c} \Proves[n] \lexists[x][A(x)], \Gamma \Sequent
  \Delta$ असेल, तर $\Log{G3c} \Proves[n] A(c), \Gamma \Sequent \Delta$,
\end{enumerate}
येथे प्रत्येक प्रसंगात~$c$ हा अनुक्रमे $\Gamma$, $\Delta$
आणि $\lforall[x][A(x)]$ किंवा $\lexists[x][A(x)]$
यांत न येणारा कोणताही स्थिरांक आहे.
\end{lem}

\begin{proof}
  आपण (1) सिद्ध करू आणि (2) प्रश्न म्हणून सोडू.
  दिलेल्या~$c$ पासून सर्व अंतर्गत आयगेन स्थिरांक
  वेगळे होतील अशी सुरक्षित नियमित नावे आधी घ्या.
  युक्तिवाद \ref{pt:seq:inv:lem:G3c-invert} च्या सिद्धतेसारखाच
  आहे. विगमनात्मक टप्प्यात पुन्हा दोन प्रसंग असतात.
  $\lforall[x][A(x)]$ प्रमुख असण्याचा प्रसंग सोपा आहे:
  शेवटच्या $\RightR{\lforall}$ निगमनाचा पूर्वपक्ष
  अपेक्षित $\Gamma
  \Sequent \Delta, A(a)$ या रूपाचा असतो आणि
  त्यावर संपणाऱ्या तत्काळ उप-सिद्धता~$\pi_1$ ची
  $\pheight{\pi_1} < \pheight{\pi}$ असते.
  शिवाय, स्वचलाच्या अटीमुळे $a$ हा
  $\Gamma$, $\Delta$, किंवा $\lforall[x][A(x)]$
  यांत येत नाही. $\pi_1$ नियमित आहे असे मानता
  येत असल्याने \ref{pt:seq:qua:lem:subst-regular} नुसार
  $\Subst{\pi_1}{c}{a}$ ही $\Gamma \Sequent \Delta, A(c)$
  ची त्याच उंची ची सिद्धता आहे.
उंची~$0$ चा प्रसंगही जपला जातो: संख्यापकीय सूत्र
अणुरूप नसल्याने समान अणुरूप सूत्राचा किंवा डावीकडील
असत्याचा साक्षी उरलेल्या संदर्भात आहे; म्हणून त्याची
जागा~$A(c)$ ने बदलल्यावरही क्रमवर्ती स्वयंसिद्धक आहे.
नव्या नियमाच्या आयगेन स्थिरांकांनाही~$c$ पासून वेगळे
नावे देऊन उरलेल्या विधानीय प्रसंगांप्रमाणे नियम पुन्हा
लावायचा.

दुसऱ्या प्रसंगात $\lforall[x][A(x)]$ हे
  शेवटच्या निगमनात प्रमुख नसते; दुसरे सूत्र प्रमुख
  असते. पुन्हा एका उदाहरणाने प्रसंग पाहू. शेवटचे
  निगमन~\RightR{\land} आहे असे समजा.
  $\Delta$ मधील एखादे सूत्र~$B \land C$
  या रूपाचे असून प्रमुख आहे. सिद्धता~$\pi$
  अशी संपते:
\[
\AxiomC{}
\RightLabel{$\pi_1$}
\Deduce$\Gamma \fCenter \Delta', \lforall[x][A(x)], B$
\AxiomC{}
\RightLabel{$\pi_2$}
\Deduce$\Gamma \fCenter \Delta', \lforall[x][A(x)], C$
\RightLabel{\RightR{\land}}
\BinaryInf$\Gamma \fCenter \Delta', \lforall[x][A(x)], B \land C$
\DisplayProof
\]
येथे $\Delta = \Delta', B \land C$.
$\Delta$ मध्ये आणि पूर्ण निष्कर्षाच्या उरलेल्या
भागातही न येणारा स्थिरांक~$c$ घ्या;
तो~$\Delta', B$ किंवा $\Delta', C$ मध्येही
येत नाही हे उघड आहे. म्हणून विगमनात्मक गृहीतक
$\pi_1$ आणि~$\pi_2$ ना लागू होते; आपल्याकडे
सिद्धता~$\pi_1'$ आणि $\pi_2'$ आहेत, ज्यांसाठी
$\pheight{\pi_1'} \le \pheight{\pi_1}$ आणि
$\pheight{\pi_2'} \le
\pheight{\pi_2}$. अपेक्षित $\Gamma \Sequent
\Delta, A(c)$ ची सिद्धता~$\pi'$ अशी:
\[
\AxiomC{}
\RightLabel{$\pi_1'$}
\Deduce$\Gamma \fCenter \Delta', A(c), B$
\AxiomC{}
\RightLabel{$\pi_2'$}
\Deduce$\Gamma \fCenter \Delta', A(c), C$
\RightLabel{\RightR{\land}}
\BinaryInf$\Gamma \fCenter \Delta', A(c), B \land C$
\DisplayProof
\]
आणि $\pheight{\pi'} = \max(\pheight{\pi_1'}, \pheight{\pi_2'})+1 \le \pheight{\pi}$.
\end{proof}

\ref{pt:seq:inv:lem:G3c-invert} चा कट-निरसन प्रमेयाच्या
सिद्धतेत महत्त्वाचा उपयोग होतो. पण त्याने पुढील
निष्कर्षही दाखवता येतो:

\begin{prop}\label{pt:seq:inv:prop:G3c-cont-adm}
\Log{G1c} चे आकुंचन नियम~\LeftR{\Contraction} आणि
\RightR{\Contraction} हे \Log{G3c} मध्ये
उंची जपून अनुज्ञेय आहेत.
\end{prop}

\begin{proof}
\RightR{\Contraction} आणि \LeftR{\Contraction}
दोन्हींसाठी एकाच वेळी युक्तिवाद देऊ.
\Log{G3c} मध्ये $\pi$ ही $\Gamma \Sequent
\Delta, A, A$ ची किंवा $A, A, \Gamma \Sequent
\Delta$ ची सिद्धता आहे असे माना. अनुक्रमे
$\Gamma \Sequent \Delta, A$ किंवा $A,
\Gamma \Sequent \Delta$ ची अशी \Log{G3c}-सिद्धता
~$\pi'$ आहे की $\pheight{\pi'} \le
\pheight{\pi}$, हे दाखवायचे आहे. ते
$n = \pheight{\pi}$ वर विगमन करून करू.

$n=0$ असेल, तर~$\pi$ हे फक्त स्वयंसिद्धक
आहे. $\Gamma \Sequent \Delta, A,
A$ हे \Log{G3c} चे स्वयंसिद्धक असेल तर
$\Gamma \Sequent \Delta, A$ हेही स्वयंसिद्धक
असते: कारण एखादा अण्वीय~$C$ हा~$\Gamma$
आणि~$\Delta$ दोहोंचा घटक असतो, किंवा
$A$ अण्वीय असून~$\Gamma$ चा घटक
असतो. अंतिम क्रमवर्ती $A, A, \Gamma \Sequent
\Delta$ असेल, तरी हाच युक्तिवाद लागू होतो.
डावीकडील असत्यामुळे स्वयंसिद्धक असेल, तर आकुंचनानंतर
किमान एक असत्य आस्थिती राहते; म्हणून तो प्रसंगही
उंची~$0$ चा असतो.

$n>0$ असेल, तर~$\pi$ चा शेवट काही
नियम~$R$ ने होतो. या शेवटच्या निगमनात $A$
प्रमुख नसेल, तर \ref{pt:seq:inv:lem:G3c-invert} च्या
सिद्धतेप्रमाणे~$\pi$ च्या तत्काळ उप-सिद्धता
ना विगमनात्मक गृहीतक लावून आणि मिळालेल्या
सिद्धता(s) ना पुन्हा तोच नियम~$R$ लावून
सिद्धता~$\pi'$ मिळवता येते. विगमनात्मक गृहीतक
असे आहे: $\Pi \Sequent \Lambda, D,
D$ किंवा $D, D, \Pi \Sequent \Lambda$ यांची
उंची~$k<n$ असलेली सिद्धता असेल, तर
अनुक्रमे $\Pi \Sequent \Lambda, D$ किंवा $D, \Pi \Sequent
\Lambda$ यांची उंची~$\le k$ असलेली
सिद्धता आहे. (संदर्भ किंवा आकुंचित होणारे
सूत्र हे~$\Gamma$, $\Delta$, किंवा~$A$
यांहून वेगळे असले तरी विगमनात्मक गृहीतक लागू होते!)

उदाहरणार्थ, $R$ हा~$\RightR{\land}$
असून~$\pi$ चा शेवट असा आहे असे समजा:
\[
\AxiomC{}
\RightLabel{$\pi_1$}
\Deduce$\Gamma \fCenter \Delta', B, A, A$
\AxiomC{}
\RightLabel{$\pi_2$}
\Deduce$\Gamma \fCenter \Delta', C, A, A$
\RightLabel{\RightR{\land}}
\BinaryInf$\Gamma \fCenter \Delta', B \land C, A, A$
\DisplayProof
\]
येथे $\Delta = \Delta', B \land C$. विगमनात्मक
गृहीतकाने अनुक्रमे $\Gamma \fCenter \Delta', B, A$
आणि $\Gamma \fCenter \Delta', C, A$ यांच्या
सिद्धता~$\pi_1'$ आणि $\pi_2'$ मिळतात, आणि
$\pheight{\pi_1'} \le  \pheight{\pi_1}$ व
$\pheight{\pi_2'} \le
\pheight{\pi_2}$ असते. मग सिद्धता~$\pi'$ अशी:
\[
\AxiomC{}
\RightLabel{$\pi_1'$}
\Deduce$\Gamma \fCenter \Delta', B, A$
\AxiomC{}
\RightLabel{$\pi_2'$}
\Deduce$\Gamma \fCenter \Delta', C, A$
\RightLabel{\RightR{\land}}
\BinaryInf$\Gamma \fCenter \Delta', B \land C, A$
\DisplayProof
\]

शेवटच्या निगमनात~$A$ प्रमुख असेल, तर
~$\pi$ च्या तत्काळ उप-सिद्धता ना व्युत्क्रमण
पूर्वप्रमेय (\ref{pt:seq:inv:lem:G3c-invert}) लावतो, मग
विगमनात्मक गृहीतक लावतो, आणि शेवटी नियम~$R$
पुन्हा लावतो. उदाहरणार्थ, $\pi$ ही
$\Gamma \Sequent \Delta, A, A$ ची सिद्धता
असून $A
\ident (B \land C)$ हे शेवटच्या निगमनात
प्रमुख आहे असे समजा. मग~$\pi$ चे रूप असे:
\[
\AxiomC{}
\RightLabel{$\pi_1$}
\Deduce$\Gamma \fCenter \Delta, B \land C, B$
\AxiomC{}
\RightLabel{$\pi_2$}
\Deduce$\Gamma \fCenter \Delta, B \land C, C$
\RightLabel{\RightR{\land}}
\BinaryInf$\Gamma \fCenter \Delta, B \land C, B \land C$
\DisplayProof
\]
\ref{pt:seq:inv:lem:G3c-invert} हे~$\pi_1$ ला लावल्यास
$\Gamma \Sequent \Delta, B, B$ ची
सिद्धता~$\pi_1'$ मिळते. (त्यातून
$\Gamma \Sequent \Delta, C, B$ ची सिद्धता
सुद्धा मिळते, पण आपल्याला ती लागत नाही.)
तसेच~$\pi_2$ ला \ref{pt:seq:inv:lem:G3c-invert} लावल्यास
$\Gamma \Sequent \Delta, C, C$ ची
सिद्धता~$\pi_2'$ मिळते. \RightR{\land} चे
व्युत्क्रमण उंची जपते, म्हणून
$\pheight{\pi_1'} \le  \pheight{\pi_1} < \pheight{\pi}$
आणि $\pheight{\pi_2'} \le  \pheight{\pi_2} < \pheight{\pi}$.
त्यामुळे~$\pi_1'$ आणि $\pi_2'$ ना विगमनात्मक गृहीतक
लागू होते: अनुक्रमे $\Gamma \Sequent \Delta, B$
आणि $\Gamma \Sequent \Delta, C$ यांच्या
सिद्धता~$\pi_1''$ आणि $\pi_2''$ मिळतात, आणि
$\pheight{\pi_1''}
\le \pheight{\pi_1'} \le \pheight{\pi_1}$ तसेच
$\pheight{\pi_2''} \le
\pheight{\pi_2'} \le \pheight{\pi_2}$ असते.
मग सिद्धता~$\pi'$ अशी:
\[
\AxiomC{}
\RightLabel{$\pi_1''$}
\Deduce$\Gamma \fCenter \Delta, B$
\AxiomC{}
\RightLabel{$\pi_2''$}
\Deduce$\Gamma \fCenter \Delta, C$
\RightLabel{\RightR{\land}}
\BinaryInf$\Gamma \fCenter \Delta, B \land C$
\DisplayProof
\]
तिच्यासाठी $\pheight{\pi'} \le \pheight{\pi}$.


सूत्र प्रमुख असताना संख्यकांच्या
नियमांबाबत विशेष काळजी घ्यावी लागते.

$A \ident \lforall[x][B(x)]$ हे उजवीकडे
प्रमुख आहे असे समजा. मग~$\pi$ चा शेवट असा:
\[
\AxiomC{}
\RightLabel{$\pi_1$}
\Deduce$\Gamma \fCenter \Delta, \lforall[x][B(x)], B(a)$
\RightLabel{\RightR{\lforall}}
\UnaryInf$\Gamma \fCenter \Delta, \lforall[x][B(x)], \lforall[x][B(x)]$
\DisplayProof
\]
\ref{pt:seq:inv:lem:invert-quant} नुसार, $\Gamma \fCenter \Delta,
\lforall[x][B(x)], B(a)$ मध्ये न येणाऱ्या
कोणत्याही~$c$ साठी $\Gamma \fCenter \Delta, B(c),
B(a)$ ची उंची~$\le \pheight{\pi_1}$
असलेली सिद्धता~$\pi_1'$ मिळते. $\pi_1'$
नियमित करून तिचे अंतर्गत आयगेन स्थिरांक~$a$ पासून
वेगळे आहेत असे मानता येते; म्हणून
\ref{pt:seq:qua:lem:subst-regular} नुसार
सिद्धता~$\Subst{\pi_1'}{a}{c}$ ही $\Gamma \fCenter \Delta, B(a),
B(a)$ ची त्याचप्रमाणे उंची~$\le \pheight{\pi_1}$
असलेली सिद्धता आहे. आता विगमनात्मक गृहीतक
लागू होऊन $\Gamma
\fCenter \Delta, B(a)$ ची सिद्धता~$\pi_1''$
मिळते. \RightR{\lforall} नियम लावून~$\pi'$
मिळवू:
\[
\AxiomC{}
\RightLabel{$\pi_1''$}
\Deduce$\Gamma \fCenter \Delta, B(a)$
\RightLabel{\RightR{\lforall}}
\UnaryInf$\Gamma \fCenter \Delta, \lforall[x][B(x)]$
\DisplayProof
\]
आणि $\pheight{\pi'} \le \pheight{\pi}$.


आता $A \ident \lexists[x][B(x)]$ हे
उजवीकडे प्रमुख आहे असे समजा. मग~$\pi$ चे रूप असे:
\[
\AxiomC{}
\RightLabel{$\pi_1$}
\Deduce$\Gamma \fCenter \Delta, \lexists[x][B(x)], \lexists[x][B(x)], B(t)$
\RightLabel{\RightR{\lexists}}
\UnaryInf$\Gamma \fCenter \Delta, \lexists[x][B(x)], \lexists[x][B(x)]$
\DisplayProof
\]
विगमनात्मक गृहीतक~$\pi_1$ ला लागू होते; त्यामुळे
$\Gamma \Sequent \Delta, \lexists[x][B(x)], B(t)$
ची सिद्धता~$\pi_1'$ मिळते.
(येथे व्युत्क्रमण पूर्वप्रमेयाची गरज नाही!) मग
सिद्धता~$\pi'$ अशी:
\[
\AxiomC{}
\RightLabel{$\pi_1'$}
\Deduce$\Gamma \fCenter \Delta,  \lexists[x][B(x)], B(t)$
\RightLabel{\RightR{\lexists}}
\UnaryInf$\Gamma \fCenter \Delta, \lexists[x][B(x)]$
\DisplayProof
\]
तिच्यासाठी $\pheight{\pi'} \le \pheight{\pi}$.
\end{proof}

\begin{prob}
\LeftR{\Contraction} साठी \ref{pt:seq:inv:prop:G3c-cont-adm}
च्या सिद्धतेची रूपरेषा द्या. $A \ident (B \land C)$
आणि $A \ident (B
\lif C)$ हे प्रसंग वर्णा.
\end{prob}

\begin{prob}
उरलेल्या संख्यकांच्या प्रसंगांसाठी, म्हणजे $A \ident
\lforall[x][B(x)]$ आणि $A \ident \lexists[x][B(x)]$
असताना, \LeftR{\Contraction} साठी
\ref{pt:seq:inv:prop:G3c-cont-adm} च्या सिद्धतेची रूपरेषा द्या.
\end{prob}













\section{\Log{G1c} आणि \Log{G3c} यांतील रूपांतरे}\label{pt:seq:tr:sec}

\Log{G1c} आणि \Log{G3c} या दोन्ही पद्धती अभिजात तर्कशास्त्रासाठी
निर्दोष व संपूर्ण आहेत, असे आपण म्हटले होते. म्हणजे प्रत्येक
संरचनांमध्ये सत्य असलेल्या आणि फक्त त्याच क्रमवर्ती त्या
सिद्ध करतात. विशेषतः, दोन्ही त्याच क्रमवर्ती सिद्ध करतात. आता
निर्दोषता व संपूर्णतेच्या प्रमेयांचा वळसा न घालता, केवळ
सिद्धता-उपपत्तीय मार्गाने हे तथ्य सिद्ध करता येते. अशा
सिद्धतेत \Log{G1c} मधील सिद्धता चे \Log{G3c} मधील
सिद्धता मध्ये आणि उलट दिशेने \emph{रूपांतरे} मिळतात.

\begin{prop}
$\Log{G1c} \Proves S$ असेल, तर $\Log{G3c} \Proves S$.
\end{prop}

\begin{proof}
  $\pi$ ही \Log{G1c} मधील सिद्धता असेल, तर \Log{G3c} मध्ये
  $S$ ची सिद्धता~$\pi'$ आहे, हे दाखवू. यासाठी $n
  = \pheight{\pi}$ वर विगमन करू.

  $n = 0$ असेल, तर $\pi$ हे स्वयंसिद्धक असते. $\lfalse \Sequent
  \quad$ हे स्वयंसिद्धक \Log{G3c} मध्येही आहे (संदर्भ $\Gamma,
  \Delta$ रिकामे घ्या). \Log{G1c} मधील $A \Sequent A$ ही
  स्वयंसिद्धके $A$ अण्वीय असण्याची अट न घालता ग्राह्य असतात.
  पण अशा सर्व क्रमवर्ती \Log{G3c} मध्ये सिद्धतायोग्य आहेत, हे आपण
  \ref{pt:seq:ex:prop:G3c-axioms} मध्ये सिद्ध केले आहे.

  $n>0$ असेल, तर $\pi$ चा शेवट $R$ नियमाने होतो. विगमनाचे
  गृहीतक सांगते की त्या नियमाच्या आधारविधानांच्या \Log{G3c}
  मधील सिद्धता आहेत; म्हणजे लगेचच्या उप-सिद्धता ची
  \Log{G3c} मध्ये रूपांतरे आहेत. $R$ हा \Log{G3c} मधीलही नियम
  असेल, तर आधारविधानांच्या रूपांतरित सिद्धता वर तो लागू करू.
  $R$ हा \Log{G1c} मधील $\LeftR{\land}$ असेल,
  तर आधारविधानात नसलेले दुसरे संयोजक सूत्र डावीकडे
  दुर्बलीकरणाने जोडा आणि \Log{G3c} चा $\LeftR{\land}$
  नियम लावा. $\RightR{\lor}$ साठी नसलेले दुसरे वियोजक
  सूत्र उजवीकडे जोडा आणि \Log{G3c} चा $\RightR{\lor}$
  नियम लावा. \Log{G1c} च्या $\LeftR{\lforall}$ किंवा
  $\RightR{\lexists}$ साठी त्याच बाजूला मुख्य संख्यापकीय
  सूत्र दुर्बलीकरणाने आधी जोडा; मग ते जपणारा \Log{G3c}
  नियम लावा. दुर्बलीकरणाची अनुज्ञेयता
  \ref{pt:seq:adm:prop:weak-G3c-adm} देते. $R$ स्वतः
  दुर्बलीकरणाचा नियम असेल, तर तीच प्रमेयिका वापरा;
  आकुंचनाचा नियम असेल, तर
  \ref{pt:seq:inv:prop:G3c-cont-adm} वापरा.
\end{proof}

\begin{prop}
$\Log{G3c} \Proves S$ असेल, तर $\Log{G1c} \Proves S$.
\end{prop}

\begin{proof}
  पुन्हा $S$ ची सिद्धता~$\pi$ हिच्या उंची वर विगमन करू.
  \Log{G3c} मधील प्रत्येक स्वयंसिद्धक \Log{G1c} मध्ये एका
  स्वयंसिद्धकापासून \LeftR{\Weakening} आणि
  \RightR{\Weakening} नियम काही वेळा वापरून सिद्ध करता येते:
  \[
  \Axiom$C \fCenter C$
  \RightLabel{\LeftR{\Weakening}}
  \Deduce$C, \Gamma \fCenter C$
  \RightLabel{\RightR{\Weakening}}
  \Deduce$C, \Gamma \fCenter \Delta, C$
  \DisplayProof
\qquad
  \Axiom$\lfalse \fCenter$
  \RightLabel{\LeftR{\Weakening}}
  \Deduce$\lfalse, \Gamma \fCenter$
  \RightLabel{\RightR{\Weakening}}
  \Deduce$\lfalse, \Gamma \fCenter \Delta$
  \DisplayProof
  \]
  $\pi$ चा शेवट $R$ नियमाने होत असेल, तर विगमनाच्या गृहीतकाने
  आधारविधानांच्या \Log{G1c} मधील सिद्धता मिळतात;
  नियम दोन्ही पद्धतींत समान असेल, तर त्या सिद्धता वर
  तोच $R$ नियम लागू करता येतो. $\LeftR{\lforall}$ आणि
  $\RightR{\lexists}$ मध्ये \Log{G3c} च्या आधारविधानात
  मुख्य संख्यापकीय सूत्रही जपलेले असते. ते संदर्भात ठेवून
  \Log{G1c} चा नियम लावल्यास निष्कर्षात मुख्य सूत्राच्या
  दोन प्रती मिळतात; अनुक्रमे डाव्या किंवा उजव्या आकुंचनाने
  त्या एक करता येतात. उरलेले अपवाद म्हणजे \LeftR{\land}
  आणि \RightR{\lor}; त्यांच्या \Log{G1c} आणि
  \Log{G3c} मधील आवृत्त्या वेगळ्या आहेत. \Log{G3c} मधील
  आवृत्त्या \Log{G1c} मध्ये पुढीलप्रमाणे व्युत्पाद्य आहेत:
  \[
  \Axiom$A, B, \Gamma \fCenter \Delta$
  \RightLabel{\LeftR{\land}}
  \UnaryInf$A \land B, B, \Gamma \fCenter \Delta$
  \RightLabel{\LeftR{\land}}
  \UnaryInf$A \land B, A \land B, \Gamma \fCenter \Delta$
  \RightLabel{\LeftR{\Contraction}}
  \UnaryInf$A \land B, \Gamma \fCenter \Delta$
    \DisplayProof
\qquad
  \Axiom$\Gamma \fCenter \Delta, A, B$
  \RightLabel{\RightR{\lor}}
  \UnaryInf$\Gamma \fCenter \Delta, A \lor B, B$
  \RightLabel{\RightR{\lor}}
  \UnaryInf$\Gamma \fCenter \Delta, A \lor B, A \lor B$
  \RightLabel{\RightR{\Contraction}}
  \UnaryInf$\Gamma \fCenter \Delta, A \lor B$
    \DisplayProof
  \]
\end{proof}
\part{इतिहास}\label{his:part}
